\documentclass{article}

% EMPIRICAL HOLD (2026-06-30): this manuscript is integrated only through SD-8.1.
% Do not treat its SD-9/landmark/sparse efficiency language as the current set-dictionary verdict or
% as launch guidance. The active evidence is exact dense and the five-seed matrix is documented in
% docs/sd_dense_paper5_matrix.md. Rewrite affected results after that matrix is validated.

% If you need to pass options to natbib, do so before loading neurips_2025.
\usepackage{neurips_2025}

\usepackage[utf8]{inputenc}
\usepackage[T1]{fontenc}
\usepackage{microtype}
\usepackage{graphicx}
\usepackage{url}
\usepackage{nicefrac}
\usepackage{xcolor}
\usepackage{subcaption}
\usepackage{booktabs}
\usepackage{longtable}
\usepackage{hyperref}
\usepackage{placeins}

\newcommand{\theHalgorithm}{\arabic{algorithm}}

\usepackage{amsmath}
\usepackage{amssymb}
\usepackage{amsfonts}
\usepackage{dsfont}
\usepackage{mathtools}
\usepackage{amsthm}
\newtheorem{framework}{Framework}
\newcommand{\mathbbm}[1]{\mathds{#1}}

% if you use cleveref..
\usepackage[capitalize,noabbrev]{cleveref}

%%%%%%%%%%%%%%%%%%%%%%%%%%%%%%%%
% THEOREMS
%%%%%%%%%%%%%%%%%%%%%%%%%%%%%%%%
\theoremstyle{plain}
\newtheorem{theorem}{Theorem}[section]
\newtheorem{proposition}[theorem]{Proposition}
\newtheorem{lemma}[theorem]{Lemma}
\newtheorem{corollary}[theorem]{Corollary}
\theoremstyle{definition}
\newtheorem{definition}[theorem]{Definition}
\newtheorem{assumption}[theorem]{Assumption}
\theoremstyle{remark}
\newtheorem{remark}[theorem]{Remark}

\usepackage[textsize=tiny]{todonotes}

\title{Set Attention}

\author{Anonymous Authors}

\begin{document}

\maketitle

\begin{abstract}
We introduce Set Attention, a sequence modeling architecture that replaces token-to-token attention with set construction, probabilistic pooling, set-level processing, and multi-head routing back to tokens. A tiered analysis separates exact structural facts, mechanistic stability bounds, and empirically testable operating-regime predictions. The theory shows that multi-head routing strictly enlarges the routing function class until limited by candidate-set topology, characterizes topology-induced entropy collapse, and decomposes gradient transport through pooling, set processing, and routing.

Controlled language-modeling experiments support these predictions across dense exact, sparse local-band, and linear landmark set-attention families. Stride/candidate-count sweeps expose topology-induced routing collapse, pooling-temperature sweeps link effective support to gradient transport, and LR-normalized matched controls show that the present causal implementation reduces long-context memory but does not outperform matched token-attention baselines in perplexity under the tested operating points. The resulting candidate-count, concentration, and transport view offers a principled basis for diagnosing set-mediated alternatives to attention and suggests a lens for related long-context problems such as cache compression, memory pruning, and learned routing.
\end{abstract}

\section{Introduction}

Attention is the dominant primitive for sequence modeling, but standard token-to-token attention ties expressive interaction to costly pairwise structure. Efficient alternatives such as sparse, local, top-$k$, low-rank, kernelized, and linear mechanisms all constrain how information is aggregated and transported. They are usually introduced as separate designs; our goal is to expose a common structural language in which those tradeoffs can be stated explicitly.

Set Attention replaces direct token-to-token interaction with a set-mediated pipeline: tokens are organized into a bank of sets, each set is pooled into a state, set states interact through a set-attention family, and a learned multi-head router maps processed set states back to tokens. In this view, information flow is governed by bank topology, pooling concentration, set-level processing, and token-to-set routing. Dense exact, sparse local-band, and linear landmark set-attention backends then become different choices inside the same architecture rather than unrelated approximations.

This formulation isolates questions that are hard to ask in ordinary attention. When does topology force routing collapse? How many distinct routing patterns can multi-head routing express? How does pooling concentration affect gradient transport? We answer these questions with a tiered analysis: Tier~A gives exact structural facts about candidate sets and routing laws; Tier~B gives local perturbation bounds under explicit smoothness assumptions; Tier~C turns these results into experimentally testable feasible-region predictions.

The main theory shows that multi-head routing strictly enlarges the routing function class over single-head routing, with intrinsic capacity growing with the number of routing heads until capped by candidate-count topology. It also characterizes routing entropy, top-$k$ constraints, topology-induced degeneracy, and a perturbation decomposition through pooling, set processing, and routing. These results motivate an operating picture in which stable behavior depends on nondegenerate candidate diversity, sufficient pooling support, and controlled gradient transport.

Our language-modeling experiments test this picture directly. Cross-family stride and pooling-temperature sweeps show that routing collapse and pooling-support effects persist across dense exact, sparse local-band, and linear landmark Set Attention. A separate LR-normalized matched comparison provides model-level context: matched token-attention baselines remain stronger in validation perplexity, while set-routed models reduce peak VRAM in the long-context slice. More broadly, the candidate-count, concentration, and transport view may apply to long-context design problems such as cache compression, memory pruning, and learned routing, where retaining state is not enough if the retained structure collapses diversity or support.

The rest of the paper is organized as follows. Section~\ref{sec:model-overview} defines the Set Attention architecture and the dense, sparse, and linear set-processing families. Section~\ref{sec:main-theory} states the compact Tier~A--C theoretical results, with full proofs deferred to the appendix. Section~\ref{sec:results} reports the Tier~C mechanism experiments and the LR-normalized baseline comparison. The appendix gives the full notation, detailed model definitions, proofs, provenance tables, and extended diagnostics.

\section{Related Work}
Prior work on structured and efficient alternatives to token-to-token attention can be organized into several overlapping themes. One line constrains token-to-token interaction directly through sparse or structured patterns such as locality, top-$k$ selection, or other restricted connectivity rules. Another line seeks linear or low-rank approximations that preserve the token-level attention interface while reducing the cost of dense pairwise interaction. These approaches are closely related to the efficiency motivation of this paper, but they typically treat the attention map itself as the main object to sparsify, approximate, or compress.

A second line of work changes how attention information is composed or aggregated rather than only changing the score kernel. LASER and DCFormer are examples within this broader design space: both explore departures from standard attention by modifying how information is organized inside the architecture, including mechanisms that alter head interaction or internal aggregation structure. We use the term \emph{routing} for our own explicit token-to-set assignment distribution; prior non-MoE transformer variants are better understood as inducing routing-like information flow rather than as implementing the same explicit routing module. The key difference is that our construction makes set formation, pooling, and routing explicit architectural stages, with bank topology entering as a first-class structural object rather than as an implicit property of an attention mask or a head-composition rule.

A third line emphasizes deployment setting or modality. Ripple Attention illustrates this perspective by tailoring attention structure to vision workloads, where the practical constraints and useful inductive biases differ from those of autoregressive language modeling. More broadly, this literature shows that efficient attention design is often shaped by hardware considerations, modality-specific structure, or the need to preserve strong baseline behavior under restricted interaction budgets.

Our paper differs from these lines of work in both representation and analysis. Instead of approximating token-to-token attention directly, we study set-mediated token interaction through explicit bank topology, probabilistic pooling, set-level processing, and routing distributions back to token space. This lets us analyze pooling concentration, candidate-count structure, and routing behavior within a single framework, and then connect that framework to a tiered theoretical framing tied to the empirical diagnostics in the Results section. That perspective also motivates the formal setup of the next section, where the bank, pooling law, and routing map are introduced as explicit mathematical objects.


%%%%%%%%%%%%%%%%%%%%%%%%%%%%%%%%%%%%%%%%%%%%%%%%%%%%%%%%%%%%%%%%%%%%%%%%%%%%%%%%
% SECTION 4: MODEL DEFINITION AND STRUCTURAL OBJECTS
%%%%%%%%%%%%%%%%%%%%%%%%%%%%%%%%%%%%%%%%%%%%%%%%%%%%%%%%%%%%%%%%%%%%%%%%%%%%%%%%


\section{Model Overview}
\label{sec:model-overview}

Set Attention replaces direct token-to-token attention by a set-mediated pipeline
\[
x_{1:L}\ \longrightarrow\ h_{1:L}^{(0)}\ \longrightarrow\ S^{(0)}\ \longrightarrow\ S^{(K)}\ \longrightarrow\ r_{1:L}\ \longrightarrow\ f_{1:L}\ \longrightarrow\ y_{1:L}.
\]
Here \(h_{1:L}^{(0)}\) are initial token states, \(S^{(0)}=(s_1^{(0)},\ldots,s_M^{(0)})\) is the matrix of pooled set states, \(S^{(K)}=(s_1^{(K)},\ldots,s_M^{(K)})\) is the matrix after \(K\) set-attention blocks, \(r_{1:L}\) are routed set-context states, \(f_{1:L}\) are the final residual token states, and \(y_{1:L}\) are vocabulary logits.
For a fixed sample, token states are
\begin{equation}
\label{eq:main-token-state}
h_t^{(0)}=e(x_t)+p_t,\qquad t\in[L],
\end{equation}
with learned token and positional embeddings. In the autoregressive strict-past setting used for the main experiments, the bank drops partial trailing windows. With starts $a_m=1+(m-1)s$ and endpoints $e_m=a_m+w-1$,
\begin{equation}
\label{eq:main-bank}
M=\left\lfloor\frac{L-w}{s}\right\rfloor+1,
\qquad
S_m=\{t\in[L]:a_m\le t\le e_m\},
\end{equation}
so every active set has exactly $w$ tokens. The strict-past routing candidate
fiber is endpoint based,
\begin{equation}
\label{eq:main-strict-past-candidates}
\mathcal C_t^{\mathrm{sp}}=\{m\in[M]:t-w<e_m\le t\}.
\end{equation}
This is the candidate family used by the causal language-modeling experiments. A legacy noncausal mode can instead use membership candidates $\{m:t\in S_m\}$, but that is not the main autoregressive paper mode.

Each bank set is pooled into a set state \(s_m^{(0)}\), the expectation of token states under a probability law $\mathbb P_m\in\Delta(S_m)$,
\begin{equation}
\label{eq:main-pooling}
s_m^{(0)}=\sum_{t\in S_m}\omega_{t,m}h_t^{(0)},
\qquad
\omega_{\cdot,m}\in\Delta(S_m).
\end{equation}
The main experiments use soft-trimmed Boltzmann pooling,
$\omega_{t,m}\propto \exp(-\delta_{t,m}/\tau_{\mathrm{pool}})r^{\mathrm{trim}}_{t,m}$,
with normalization over $S_m$; mean pooling is the corresponding uniform law. The full definition is in Appendix~\ref{sec:pooling-operator} and uses $|S_m|$ wherever set cardinality matters.

The set states are processed by $K$ residual set-attention blocks with set-attention heads distinct from routing heads. For one set-attention head,
\begin{equation}
\label{eq:main-set-score}
a_{ij}^{(h)}=
\frac{\langle q_i^{(h)},k_j^{(h)}\rangle}{\sqrt{d_a}}
+g_{ij}+b_{ij}^{(h)}+m_{ij},
\end{equation}
where $q_i^{(h)},k_j^{(h)}\in\mathbb R^{d_a}$ are set-level query/key projections, $g_{ij}$ is geometry bias, $b_{ij}^{(h)}$ is learned content bias, and $m_{ij}$ is the causal or structural mask term. These additional structural terms are part of the dense exact set-attention backend. Let $A^{(h)}=(a_{ij}^{(h)})$ denote the resulting score matrix, and let $Q,K,V$ denote the corresponding query, key, and value matrices for that set-attention head.

The set-processing families used in this work are:
\begin{align}
\operatorname{SetDense}(Q,K,V) &= \operatorname{softmax}(A)V,\\
\operatorname{SetLinear}(Q,K,V) &= \operatorname{softmax}(A_{M\Lambda})\operatorname{softmax}(A_{\Lambda M})V,\\
\operatorname{SetSparse}(Q,K,V) &: A_{ij}=-\infty\quad\text{outside a local band.}
\end{align}
The linear family is the landmark backend, where $\Lambda$ indexes landmark set states, and the sparse family is the local-band backend.

Routing maps processed set states back to token states. With descriptor $d_m$ for set $m$, the learned multi-head router computes
\begin{equation}
\label{eq:main-router-score}
q_t^{(u)}=W_{Q,r}^{(u)}h_t^{(0)},\qquad
k_m^{(u)}=W_{K,r}^{(u)}d_m,
\qquad
g_{t,m}^{(u)}=\frac{\langle q_t^{(u)},k_m^{(u)}\rangle}{\sqrt{d_\phi}}.
\end{equation}
In the hashed-count experimental path, the router descriptor is the pooled set state, $d_m=s_m^{(0)}$; hashed-count features enter the set-attention content bias. Candidate masking sets non-candidates to $-\infty$; the optional fixed top-$k$ step also sets all but the retained candidates to $-\infty$. Denote the resulting masked logit by $\tilde g_{t,m}^{(u)}$. Routing probabilities are
\begin{equation}
\label{eq:main-routing-probs}
\pi_{t,m}^{(u)}\propto \exp\!\left(\tilde g_{t,m}^{(u)}/\tau_r^{\mathrm{eff}}\right),
\qquad m\in\mathcal C_t,
\end{equation}
where $\tau_r^{\mathrm{eff}}=\max\{\tau_r,\tau_{\min}\}$ for the configured lower clamp $\tau_{\min}>0$. Writing $\tilde s_m^{(u)}$ for the head-$u$ subvector of the processed set state $s_m^{(K)}$, the routed head output is
\begin{equation}
\label{eq:main-routed-head}
r_t^{(u)}=\sum_{m\in\mathcal C_t}\pi_{t,m}^{(u)}\tilde s_m^{(u)},
\qquad
r_t=\operatorname{Concat}_{u=1}^{H_r}r_t^{(u)}.
\end{equation}
If $\mathcal C_t=\varnothing$, the routed context is defined as $r_t=0$.
The default language-model head receives the direct residual representation
\begin{equation}
\label{eq:main-direct-residual}
f_t=h_t^{(0)}+r_t.
\end{equation}
For calibration experiments that study convergence toward the token-level
limit, we also use a named output policy
\begin{equation}
\label{eq:main-empty-only-residual}
f_t^{\mathrm{empty}}=
\begin{cases}
h_t^{(0)},& C_t=0,\\
r_t,& C_t>0.
\end{cases}
\end{equation}
This \texttt{empty\_only} policy preserves a well-defined state for early
strict-past positions with no sealed endpoint while avoiding an additional
token residual whenever routed set information is present. The \texttt{direct}
policy in \eqref{eq:main-direct-residual} is the default for the reported
A1--A6 runs; \texttt{empty\_only} is used only when explicitly configured for
token-limit and compression-calibration experiments.

A third output policy, the \texttt{anchor\_span} residual analyzed in
Appendix~\ref{sec:future-set-dictionary}, reuses the same final-state form as
\eqref{eq:main-direct-residual},
\begin{equation}
\label{eq:main-anchor-span}
f_t = h_t^{(0)} + r_t,
\end{equation}
but disables the optional pre-pooling token MLP (Appendix~\ref{app:token-mlp}) so
that $h_t^{(0)}=e(x_t)+p_t$ is a deliberately thin identity carrier and the routed
context $r_t$ of \eqref{eq:main-routed-head} must supply the contextual prediction.
A span ablation that sets $r_t{:=}0$ at evaluation collapses perplexity toward the
unigram regime, confirming that $h_t^{(0)}$ is not a predictive token bypass and
that prediction is genuinely set-mediated. Equation~\eqref{eq:set-dictionary-span}
writes $r_t$ through the set-output projection $W_o$.\footnote{$W_o$ maps the
concatenated routed context from the set-state width to $D$; it is the identity
whenever the two widths coincide, as in every reported configuration, so it is
omitted from \eqref{eq:main-routed-head}.}

\paragraph{Multi-resolution head groups.}
The $H_r$ routing heads may be partitioned into $G$ disjoint groups indexed by
$\nu$, with head sets $\mathcal U_\nu\subseteq\{1,\dots,H_r\}$ satisfying
$\bigsqcup_{\nu}\mathcal U_\nu=\{1,\dots,H_r\}$. Each group $\nu$ carries its own
bank $\mathcal H_{L,w_\nu,s_\nu}$, and hence its own sets $S^{(\nu)}_m$, endpoints
$e^{(\nu)}_m$, candidate fiber
$\mathcal C^{(\nu)}_t=\{m:t-w_\nu<e^{(\nu)}_m\le t\}$ (mean count
$\bar C^{(\nu)}\approx w_\nu/s_\nu$), pooled states $s^{(0,\nu)}_m$ by
\eqref{eq:main-pooling}, and processed atoms $s^{(K,\nu)}_m$ from a separate set
stack. A head $u\in\mathcal U_\nu$ routes over \emph{its own} fiber
$\mathcal C^{(\nu)}_t$ in \eqref{eq:main-routing-probs}--\eqref{eq:main-routed-head},
and the head subvectors $r^{(u)}_t$ are concatenated across all groups to form
$r_t$. We use a fine group $(w_{\mathrm f},s_{\mathrm f})=(2,1)$ and a coarse,
``blurred'' group $(w_{\mathrm c},s_{\mathrm c})=(4,2)$; the coarse-head fraction
$|\mathcal U_{\mathrm c}|/H_r$ sets the resolution mix, with $G=1$ the
single-resolution model and the fractions $0$ and $1$ recovering the uniform fine
and coarse banks.

Finally, vocabulary logits satisfy $y_t\in\mathbb R^{V_{\mathrm{vocab}}}$. Full notation, diagnostics, and optional implementation components are given in Appendix~\ref{sec:model-objects} and Appendix~\ref{app:token-mlp}.



\section{Theory}
\label{sec:main-theory}

The formal analysis below is for the executed multiresolution set-dictionary branch: exact dense set attention, separate fine and coarse banks, endpoint-window candidate gather, token MLP disabled, trained anchor disabled, and \texttt{anchor\_span} readout.  We state four compact results here and give complete proofs, code correspondence, and legacy-label disposition in Appendix~\ref{sec:theory}.  These statements are structural or conditional; they do not assert a perplexity win, token-attention equivalence, subquadratic complexity, or an intrinsic OOM theorem.

\begin{theorem}[Contextual-path factorization]
\label{thm:main-contextual-path-factorization}
For two histories $x,x'$ with the same current token and position at $t$, the active \texttt{anchor\_span} model satisfies
\begin{equation}
\label{eq:main-contextual-path-factorization}
\ell_t(x)-\ell_t(x')=W_{\rm lm}W_o\bigl(r_t(x)-r_t(x')\bigr).
\end{equation}
Thus every dependence of the logits on $x_{<t}$ factors through the routed span $W_or_t$.  Span differences in the kernel of $W_{\rm lm}W_o$ up to an all-classes logit shift are not observable from softmax probabilities alone.
\end{theorem}

\begin{theorem}[Direct-sum multigroup routing capacity]
\label{thm:main-multigroup-routing-capacity}
Let $\mathcal G$ be the set of resolution groups, $\mathcal U_g$ the routing heads in group $g$, $H_g=|\mathcal U_g|$, and $\mathcal C_t^g$ the group endpoint fiber.  If $\widehat\Pi_t$ is the routing matrix over the tagged disjoint union $\bigsqcup_g(\{g\}\times\mathcal C_t^g)$, then
\begin{equation}
\label{eq:main-multigroup-routing-capacity}
\operatorname{rank}(\widehat\Pi_t)=\sum_g \operatorname{rank}(\Pi_t^g)
\le \sum_g \min\{H_g,|\mathcal C_t^g|\}.
\end{equation}
The admissible assignment space is the product of group-local simplices.  Finite-temperature dot-product routers parameterize smooth subsets of this space; boundary Dirac laws are limiting cases, not generally finite-temperature realizations.
\end{theorem}

\begin{theorem}[Exact-dense score-memory coefficient]
\label{thm:main-exact-dense-memory-coefficient}
For sequence length $L$, batch size $B$, $K$ set-attention blocks, and group bank sizes $M_g(L)=\lfloor(L-w_g)/s_g\rfloor+1$, the active exact backend materializes
\begin{equation}
\label{eq:main-score-count}
A_{\rm score}(L,B)=BK\sum_g H_gM_g(L)^2
\end{equation}
score elements before masking.  For fine $(w,s)=(2,1)$ and coarse $(4,2)$ with coarse-head fraction $p$,
\begin{equation}
\label{eq:main-blur-coefficient}
\frac{A_{\rm score}}{BKHL^2}=1-\frac{3p}{4}+O(L^{-1}).
\end{equation}
This is a constant-factor change inside the exact-dense quadratic family.
\end{theorem}

\begin{theorem}[Conditional interior allocation]
\label{thm:main-conditional-interior-allocation}
Let $n$ be the number of coarse heads, $A(n)=(H-n)M_f^2+nM_c^2$, and $E(n)$ an abstract validation or approximation error under a fixed control tuple.  If $M_c<M_f$, then $A(n)$ strictly decreases with $n$.  If the marginal gains $G(n)=E(n-1)-E(n)$ are nonincreasing with $G(1)>0$ and $G(H)<0$, then $E$ has an interior minimizer.  A mixed allocation is Pareto-better than all-fine exactly when it has no larger error and strictly smaller score memory.  The premises are conditional and must be established empirically for any specific task.
\end{theorem}
\section{Results}
\label{sec:results}

\paragraph{Protocol.}
All language-modeling results use WikiText-2 with the LR-normalized reference
width $D=384$, $d_{\mathrm{ff}}=1536$, window size $w=16$, stride $s=8$,
and learning rate $10^{-4}$ unless otherwise stated.  For set-only models,
the autoregressive setting uses strict-past endpoint routing, drops partial
trailing windows, and therefore has $M=63$ sets at $L=512$ and $M=255$ sets at
$L=2048$.  Linear-family rows use the landmark backend with coverage $0.25$.
All reported means are over three seeds unless explicitly marked otherwise;
artifact hashes and per-run metadata are recorded in
\texttt{audit/A5\_4\_handoff.md}.

\subsection{Matched Backend Controls at the LR-Normalized Reference}
\label{sec:results-matched-controls}

Table~\ref{tab:matched-headline} compares each SKA family against the matched
token-attention backend at the $L=512$ LR-normalized headline setting.  The
matched controls are essential: comparing all SKA variants only to a dense
token baseline would conflate the effect of set routing with the effect of the
underlying dense, sparse, or landmark backend.

\begin{table}[t]
\centering
\caption{Matched backend controls at $L=512$, $w=16$, $s=8$, $M=63$.  Values
are mean validation perplexity over three seeds; lower is better.  Source
artifacts: \texttt{a6\_set\_state\_width\_summary.tsv} and
\texttt{a2\_baseline\_controls\_summary.tsv}.}
\label{tab:matched-headline}
\label{tab:matched-baseline-vs-setonly}
\scriptsize
\begin{tabular}{llcc}
\toprule
Family & Backend & Token PPL $\downarrow$ & SKA PPL $\downarrow$ \\
\midrule
Dense & exact & $781.1 \pm 37.1$ & $1422.8 \pm 45.8$ \\
Sparse & local band & $740.2 \pm 3.8$ & $1527.5 \pm 47.1$ \\
Linear & landmark & $973.9 \pm 13.7$ & $1515.9 \pm 135.3$ \\
\bottomrule
\end{tabular}
\end{table}

The matched comparison is unfavorable to SKA at this budget: all three token
baselines achieve lower perplexity than their set-routed counterparts.  This
does not invalidate the structural claims, but it narrows the empirical claim:
the present implementation is a controlled diagnostic model rather than a
drop-in perplexity improvement over matched token attention.

\begin{figure*}[t]
\centering
\includegraphics[width=\textwidth]{../plots/main/fig_causal_lr_sweep_dispersion.png}
\caption{Post-A1 causal LM LR sweep at the $D=384$, $d_{\mathrm{ff}}=1536$,
$L=512$, $w=16$, $s=8$ reference.  Token baselines and SKA variants are shown
with matched dense, sparse local-band, and linear landmark backends.  Lines show
three-seed means, shaded regions show 95\% confidence intervals, and faint
markers show individual seeds.  Across the swept learning rates, matched token
baselines remain lower in perplexity while SKA variants use less memory.}
\label{fig:causal-lr-sweep-dispersion}
\end{figure*}

Figure~\ref{fig:causal-lr-sweep-dispersion} gives the corresponding
multi-seed LR view rather than only the best-row table.  The dispersion is
important for interpretation: SKA variants keep a consistent memory advantage
at this reference size, but the validation-perplexity gap persists throughout
the LR grid rather than being explained by a single unlucky learning rate.

\subsection{Calibrated Token-Limit Compression Path}
\label{sec:results-empty-only-calibration}

The matched-control table evaluates the default direct residual architecture at
the compressed LR-normalized operating point.  To test the convergence
hypothesis more directly, we also run a calibrated \texttt{empty\_only}
SetDense path.  In this mode the token residual is used only for positions with
no sealed endpoint ($C_t=0$); otherwise the language-model head receives the
routed set context.  Proposition~\ref{prop:empty-strict-past-fibers} shows that
the singleton setting $w=1,s=1$ has $M=L$ and no empty-fiber positions, making
it the cleanest empirical token-limit calibration available in the current
implementation.

\begin{figure*}[t]
\centering
\includegraphics[width=\textwidth]{../plots/main/fig_a7_empty_only_calibration.png}
\caption{Calibrated \texttt{empty\_only} SetDense compression path at
$D=384$, $d_{\mathrm{ff}}=1536$, $L=512$, and learning rate $10^{-4}$.
Validation perplexity is plotted against set-state ratio $M/L$, compression
factor $L/M$, and measured mean candidate count; peak VRAM gives the memory
tradeoff.  Three-seed SetDense points and the reused dense token baseline show
95\% confidence intervals.
The $w=1,s=1,M=L$ set point nearly reaches the dense token baseline but remains
slightly worse, while stronger compression quickly increases perplexity.}
\label{fig:a7-empty-only-calibration}
\end{figure*}

Figure~\ref{fig:a7-empty-only-calibration} separates the residual calibration
question from the compressed default setting.  The $M/L=1$ point reaches mean
validation perplexity $800.8$ versus $781.1$ for the matched dense token
baseline, using the same seed set $\{0,1,2\}$ and the same training budget.
Thus the calibrated set-mediated path approaches, but does not equal, token
attention.  As $M/L$ decreases, validation perplexity rises and peak VRAM
falls, matching the intended compression tradeoff.  This supports an empirical
convergence picture rather than an exact-equivalence claim: even in the
singleton limit, the computation still passes through set pooling, set-stack
processing, and routing projections.

\begin{figure*}[t]
\centering
\includegraphics[width=\textwidth]{../plots/main/fig_a7_backend_family_compression.png}
\caption{Backend-matched calibrated \texttt{empty\_only} compression paths for
dense exact, sparse local-band, and linear landmark families at
$D=384$, $d_{\mathrm{ff}}=1536$, $L=512$, and learning rate $10^{-4}$.
Each column compares a set family against its matched token-attention backend
baseline.  The convergence-critical singleton and near-singleton rows use five
seeds; the remaining compressed points use the original three-seed grid.  Rows
show the same validation perplexity estimates as a function of set-state ratio
$M/L$, compression factor $L/M$, and mean candidate count.  Token baselines
appear as horizontal references because they do not have a set-router candidate
fiber or a set topology.  Near $M/L=1$, the set variants approach the matched
token baselines, with the linear singleton point overlapping the landmark-token
band.  Under compression, validation perplexity rises sharply across all
backends, so the evidence supports empirical calibration rather than an
architecture-independent equivalence.}
\label{fig:a7-backend-family-compression}
\end{figure*}

The backend-family extension in
Figure~\ref{fig:a7-backend-family-compression} answers whether the dense
calibration behavior is an artifact of exact attention.  It is not: at
$w=1,s=1,M=L$, SetDense reaches mean validation perplexity $805.1$ versus
$781.1$ for the matched dense token baseline, SetSparse reaches $747.9$ versus
$742.4$ for the matched sparse token baseline, and SetLinear reaches $927.0$
versus $956.2$ for the matched landmark-token baseline, all with overlapping
five-seed confidence intervals.  The agreement does not persist under
meaningful compression: at $(w,s)=(2,1)$ the result is mixed across backends,
and by $(w,s)=(3,1)$ all three set families separate from their token
references.  Thus backend matching strengthens the operating-region
interpretation: calibrated set attention can approach matched token attention
near the uncompressed limit, but the present implementation pays
a substantial quality cost as soon as the set bank meaningfully compresses the
sequence.

\begin{figure*}[t]
\centering
\includegraphics[width=\textwidth]{../plots/main/fig_a7_seed_extension_fine_grained.png}
\caption{A7 quality--efficiency frontier for backend-matched calibrated
\texttt{empty\_only} runs.  Each panel plots the set-minus-token perplexity
gap against peak VRAM as a percentage of the matched token backend.  The star
marks the token baseline at $(100\%,0)$, and lower-left is better.  Each set
point is labeled by topology $(w,s)$ and compression factor $L/M$; error bars
show 95\% confidence intervals using three to five seeds.  Compressed points
move left, showing a memory benefit, but generally move far upward in
perplexity.  The singleton points remain near the token baselines but use
slightly more VRAM.}
\label{fig:a7-seed-extension-fine-grained}
\end{figure*}

\begin{table*}[t]
\centering
\small
\caption{Compact A7 operating-point comparison.  The top block reports
matched token-backend baselines relative to the dense token baseline.  The
bottom block reports calibrated set \texttt{empty\_only} operating points
relative to their matched token backend.  Near-token rows use five seeds;
the first meaningful compression row near $L/M\approx2$ uses the original
three-seed A7 grid.}
\label{tab:a7-operating-point-comparison}
\resizebox{\textwidth}{!}{%
\begin{tabular}{llrrrrrl}
\toprule
Group & Backend / family & Topology & $L/M$ & Val. PPL & $\Delta$PPL & VRAM & $n$ \\
\midrule
\multicolumn{8}{l}{\textit{Token baselines: $\Delta$PPL is relative to dense token attention; VRAM is absolute peak usage.}} \\
Token & Dense exact & -- & -- & 781.1 & +0.0 & 13.1 GiB & 5 \\
Token & Sparse local-band & -- & -- & 742.4 & -38.7 & 13.1 GiB & 5 \\
Token & Linear landmark & -- & -- & 956.2 & +175.1 & 12.3 GiB & 5 \\
\midrule
\multicolumn{8}{l}{\textit{Set \texttt{empty\_only}: $\Delta$PPL and VRAM are relative to the matched token backend.}} \\
Set & Dense exact & $(1,1)$ singleton & 1.000 & 805.1 & +24.0 & 103.9\% & 5 \\
Set & Dense exact & $(2,1)$ near-token & 1.002 & 867.4 & +86.3 & 103.9\% & 5 \\
Set & Dense exact & $(3,1)$ near-token & 1.004 & 922.1 & +141.0 & 103.8\% & 5 \\
Set & Dense exact & $(4,2)$ first compression & 2.008 & 1273.6 & +492.5 & 88.1\% & 3 \\
\addlinespace
Set & Sparse local-band & $(1,1)$ singleton & 1.000 & 747.9 & +5.6 & 103.9\% & 5 \\
Set & Sparse local-band & $(2,1)$ near-token & 1.002 & 810.1 & +67.7 & 103.9\% & 5 \\
Set & Sparse local-band & $(3,1)$ near-token & 1.004 & 880.6 & +138.2 & 103.8\% & 5 \\
Set & Sparse local-band & $(4,2)$ first compression & 2.008 & 1438.5 & +696.1 & 88.1\% & 3 \\
\addlinespace
Set & Linear landmark & $(1,1)$ singleton & 1.000 & 927.0 & -29.2 & 104.2\% & 5 \\
Set & Linear landmark & $(2,1)$ near-token & 1.002 & 997.4 & +41.2 & 104.2\% & 5 \\
Set & Linear landmark & $(3,1)$ near-token & 1.004 & 1114.2 & +158.0 & 104.1\% & 5 \\
Set & Linear landmark & $(4,2)$ first compression & 2.008 & 1356.2 & +399.9 & 92.3\% & 3 \\
\bottomrule
\end{tabular}%
}
\end{table*}

Figure~\ref{fig:a7-seed-extension-fine-grained} reframes the same evidence as a
quality--efficiency frontier.  With five seeds, the singleton topology has
small or overlapping deltas in all backend families ($+24.0$ PPL for dense,
$+5.6$ for sparse, and $-29.2$ for landmark linear), but these points sit at
roughly $104\%$ of the matched token VRAM rather than below it.  A focused
VRAM audit attributes this overhead to the active architecture and
implementation path---the simultaneous token and set streams together with
dense token-to-set routing and set-attention tensors---rather than to a
removable diagnostic allocation.  The compressed
points expose the only clear resource benefit: they reduce peak VRAM by roughly
$12$--$20\%$ depending on backend and topology.  Table~\ref{tab:a7-operating-point-comparison}
shows why this is not yet an acceptable quality-efficiency trade: the first
meaningful compression rows have $+399.9$ to $+696.1$ PPL gaps, much larger
than the $+175.1$ PPL gap between linear and dense token baselines in this
run set.  Thus the present evidence supports two separate conclusions: the
calibrated singleton limit is the clean quality-convergence point, while
compression can buy memory but not yet at an acceptable quality cost.

\subsection{Mechanism Sweeps}
\label{sec:results-mechanism-sweeps}

The fixed-stride window sweep and pooling-temperature sweep remain useful as
mechanism checks, but they should be interpreted separately from the headline
quality comparison.  At fixed stride $s=4$, increasing the window size changes
candidate-count structure under the post-A1 strict-past bank definition
$M=\lfloor(512-w)/4\rfloor+1$.  The summaries in
Appendix~\ref{app:complete-experiment-record} report the active dense,
sparse, and landmark set-family window and pooling-temperature sweeps with their
per-run provenance.  Appendix~\ref{app:legacy-evidence-exclusion} states the
legacy-artifact exclusion rule.
Figures~\ref{fig:causal-window-diagnostics-dispersion}
and~\ref{fig:causal-pooltau-diagnostics-dispersion} restore the corresponding
post-A1 causal LM diagnostic views with seed dispersion.  Token baselines appear
where the metric is defined; routing and pooling diagnostics are set-family
quantities.

\begin{figure*}[t]
\centering
\includegraphics[width=\textwidth]{../plots/main/fig_causal_window_diagnostics_dispersion.png}
\caption{Fixed-stride window sweep under strict-past routing at $s=4$.
Validation perplexity and peak memory include matched sparse and linear token
baselines alongside dense, sparse, and linear SKA families.  Routing entropy
and top-1 concentration are shown only for SKA rows because token-attention
baselines do not have a set-router candidate fiber.  Lines show three-seed
means, shaded regions show 95\% confidence intervals, and faint markers show
individual seeds.}
\label{fig:causal-window-diagnostics-dispersion}
\end{figure*}

\begin{figure*}[t]
\centering
\includegraphics[width=\textwidth]{../plots/main/fig_causal_pooltau_diagnostics_dispersion.png}
\caption{Pooling-temperature diagnostics for post-A1 causal LM SKA families at
$w=16$, $s=4$.  Increasing $\tau_{\mathrm{pool}}$ broadens effective pooling
support toward the full window; dense and sparse SKA continue improving through
$\tau_{\mathrm{pool}}=0.95$, while the landmark linear backend is best near
$\tau_{\mathrm{pool}}=0.5$ and degrades at $0.95$.  Lines show three-seed means,
shaded regions show 95\% confidence intervals, and faint markers show
individual seeds.}
\label{fig:causal-pooltau-diagnostics-dispersion}
\end{figure*}

These figures use only the post-A1 causal LM manifests and matched-control
artifacts.  They support the Tier~C diagnostic picture: candidate topology
controls routing concentration, broader pooling support improves transport, and
neither effect by itself closes the matched-token perplexity gap.

\subsection{Long-Context Memory and Quality}
\label{sec:results-long-context}

At $L=2048$, strict-past set-only models use $M=255$ sealed sets.  The long
context slice in Table~\ref{tab:long-context} shows the main empirical
tradeoff and includes the matched $L=512$ reference point.  SKA variants have
substantially lower peak VRAM at long context, but they still have worse
validation perplexity.  The matched sparse token baseline is both
memory-heavier and much stronger in perplexity than the corresponding sparse
SKA row.

\begin{table*}[t]
\centering
\caption{Long-context slice with matched short-context reference.  Short-context
rows use $L=512$ and, for SKA, $w=16$, $s=8$, $M=63$; long-context rows use
$L=2048$, $w=16$, $s=8$, $M=255$.  Values are mean validation perplexity over
three seeds.  Gap columns report SKA minus its matched token backend, so a
negative gap change means that the longer context narrows the SKA--token gap.
Long-context VRAM is shown as a percentage of the matched token backend, with
MiB in parentheses.  Source artifacts: \texttt{a2\_grid\_all\_runs.tsv},
\texttt{a2\_baseline\_controls\_all\_runs.tsv},
\texttt{a42\_slice\_all\_runs.tsv}, and
\texttt{a4\_long\_context\_baseline\_controls\_all\_runs.tsv}.}
\label{tab:long-context}
\scriptsize
\resizebox{\textwidth}{!}{%
\begin{tabular}{llccrrrr}
\toprule
Family & Backend & $L=512$ PPL $\downarrow$ & $L=2048$ PPL $\downarrow$ &
Gap@512 & Gap@2048 & Gap change & $L=2048$ VRAM \\
\midrule
Token dense & exact & $781.1 \pm 45.5$ & $928.2 \pm 50.2$ & -- & -- & -- & $100.0\%$ (18620) \\
Token sparse & local band & $740.2 \pm 4.7$ & $814.2 \pm 22.9$ & -- & -- & -- & $100.0\%$ (18644) \\
Token linear & landmark & $973.9 \pm 16.8$ & $984.7 \pm 23.7$ & -- & -- & -- & $100.0\%$ (15207) \\
\midrule
SKA dense & exact & $1486.8 \pm 50.5$ & $1511.5 \pm 110.8$ & $+705.7$ & $+583.4$ & $-122.3$ & $59.3\%$ (11043) \\
SKA sparse & local band & $1477.4 \pm 49.2$ & $1429.4 \pm 33.5$ & $+737.3$ & $+615.2$ & $-122.1$ & $59.2\%$ (11043) \\
SKA linear & landmark & $1500.4 \pm 26.0$ & $1434.7 \pm 37.9$ & $+526.5$ & $+450.0$ & $-76.5$ & $72.3\%$ (10996) \\
\bottomrule
\end{tabular}%
}
\end{table*}

The short-context comparison adds one useful signal.  Increasing context from
$512$ to $2048$ does not close the gap, but it narrows the SKA--token gap by
$76$--$122$ PPL while preserving a substantial long-context memory reduction
($27.7$--$40.8\%$ relative to the matched token backend).  This makes the
long-context setting a more favorable operating regime than the short
causal-LM benchmark, but the absolute perplexity gaps remain too large for a
performance-dominance claim.

\subsection{Convergence Check}
\label{sec:results-convergence}

The 30-epoch convergence panel confirms that the 10-epoch gap is not simply a
short-training artifact.  At seed 0 and the best available learning rate from
the LR-normalized summaries, token baselines remain ahead after 30 epochs:
dense token, sparse token, and linear token reach validation perplexities
1140.9, 1116.2, and 1331.9, respectively, while the dense, sparse, and linear
SKA rows reach 2439.7, 2251.1, and 2138.7.  These runs are intentionally
diagnostic and should not be written as convergence wins for SKA.

\subsection{Capacity Diagnostics}
\label{sec:results-capacity}

The A6 ablations test whether the quality gap is explained by obvious capacity
bottlenecks.  Table~\ref{tab:capacity-diagnostics} summarizes the result.  The
main conclusion is negative: moderate interface widening helps some rows, but
neither wider set states nor deeper set stacks provide a broad remedy.

\begin{table*}[t]
\centering
\caption{Capacity diagnostics at $L=512$, $w=16$, $s=8$, and LR $10^{-4}$.
Entries are mean validation perplexity over three seeds.  Source artifacts:
\texttt{a6\_dphi\_capacity\_summary.tsv},
\texttt{a6\_set\_state\_width\_summary.tsv},
\texttt{a6\_interface\_bottleneck\_summary.tsv}, and
\texttt{a64\_depth\_sweep\_summary.tsv}.}
\label{tab:capacity-diagnostics}
\scriptsize
\resizebox{\textwidth}{!}{%
\begin{tabular}{llcccc}
\toprule
Test & Family & Setting A & PPL A $\downarrow$ & Setting B & PPL B $\downarrow$ \\
\midrule
$d_\phi$ sweep & Set Dense & $d_\phi=384$ & 1422.8 & $d_\phi=768$ & 1541.0 \\
$d_\phi$ sweep & Set Sparse & $d_\phi=384$ & 1527.5 & $d_\phi=512$ & 1447.3 \\
$d_\phi$ sweep & Set Linear & $d_\phi=384$ & 1515.9 & $d_\phi=512$ & 1484.1 \\
set-state width & Set Dense & $d_s=384$ & 1422.8 & $d_s=768$ & 1520.8 \\
set-state width & Set Sparse & $d_s=384$ & 1527.5 & $d_s=512$ & 1407.6 \\
set-state width & Set Linear & $d_s=384$ & 1515.9 & $d_s=768$ & 1554.8 \\
interface match & Set Dense & $(d_s,d_\phi)=(768,384)$ & 1520.8 & $(768,768)$ & 1556.8 \\
interface match & Set Sparse & $(d_s,d_\phi)=(512,384)$ & 1407.6 & $(512,512)$ & 1421.9 \\
interface match & Set Linear & $(d_s,d_\phi)=(768,384)$ & 1554.8 & $(768,512)$ & 1427.5 \\
depth sweep & Set Dense & 6 layers & 1422.8 & 10 layers & 1568.8 \\
depth sweep & Set Sparse & 6 layers & 1527.5 & 10 layers & 1594.3 \\
depth sweep & Set Linear & 6 layers & 1515.9 & 10 layers & 1582.2 \\
\bottomrule
\end{tabular}%
}
\end{table*}

These results support a more nuanced diagnosis.  The routing interface width
$d_\phi$ can be a bottleneck in some cases, especially for the landmark family,
where $(d_s,d_\phi)=(768,512)$ improves over $(768,384)$.  However, matching
$d_\phi$ to the full set-state width often worsens perplexity, and increasing
set-stack depth from 6 to 8 or 10 layers consistently hurts validation
perplexity.  The current SKA gap is therefore not explained by a single missing
capacity knob; it is more likely a combination of routing optimization,
generalization, and interface conditioning.

\section{Conclusion}
The revised experiments establish a conservative empirical picture.  The
strict-past SKA implementation is causal, configurable, and memory-efficient at
long context, but matched token-attention controls remain stronger on
perplexity.  The capacity ablations rule out the simplest explanations that
the gap is merely due to too-small $d_\phi$, too-small set-state width, too
shallow set processing, insufficient anchoring signal, or too narrow a causal
candidate fiber.  The most defensible contribution is therefore a causal
set-mediated architecture and diagnostic analysis of when its topology,
pooling, interface width, and set-processing depth help or fail, rather than a
claim of perplexity dominance over matched token-attention baselines.

\FloatBarrier

\section{Future Work}

\subsection{Routing Memory and Long-Context Priority}
\label{sec:future-routing-memory}

The VRAM audit changes the priority list.  The current problem is not mainly
extra memory from logging or diagnostics: the near-token overhead comes from the
active architecture and implementation path, especially the simultaneous token
and set streams, dense token-to-set routing tensors, and set-attention tensors.
Thus memory should be reduced by changing the active computation, not by
subtracting a measurement artifact from the reported values.  The clearest
implementation target is candidate-gather routing.  Under strict-past
Option--1, token $t$ only routes over its endpoint fiber $\mathcal C_t$, but the
current learned router still constructs dense $[B,H,L,M]$ score and probability
tensors.  A semantics-preserving implementation should gather the active
candidates into $[B,H,L,C_{\max}]$, where $C_{\max}$ is the maximum candidate
count in the batch, and scatter the routed context back to tokens.  This directly
attacks the audit finding while preserving the causal bank definition.  A second
implementation target is to use fused or memory-efficient kernels for routing
and set self-attention where the masks allow it, following the IO-aware lesson of
FlashAttention \citep{dao2022flashattention}.  This may reduce peak memory
without changing the mathematical model, although it does not by itself solve
the quality gap.

The more important scientific direction is to test settings that should favor a
set-mediated model.  The long-context comparison in Table~\ref{tab:long-context}
is encouraging in this limited sense: the SKA--token gap narrows at $L=2048$
while SKA keeps a large VRAM reduction.  This is consistent with the broader
long-sequence literature, where the value of alternative attention mechanisms is
clearest under long-context memory pressure, sparse or routed computation, or
compressed memory rather than under short-context token-equivalent replacement
\citep{dai2019transformerxl,rae2019compressive,beltagy2020longformer,zaheer2020bigbird,roy2020routing}.  It is also aligned with latent and set bottleneck
architectures, which are usually trained around a compression interface rather
than evaluated only as a drop-in token-attention substitute
\citep{lee2019settransformer,jaegle2021perceiverio}.

\subsection{Set-Dictionary Span as a Diagnostic Branch}
\label{sec:future-set-dictionary}

The set-dictionary branch isolates a different representation path from the
default pooling--routing model.  The default strict-past model in
Section~\ref{sec:model-overview} routes a set context $r_t$ and then uses a
token residual, for example $f_t=h_t^{(0)}+r_t$ in the direct-residual policy.
The set-dictionary variant removes that contextual token residual from the
prediction path.  It keeps only the thin identity anchor $h_t^{(0)}=e(x_t)+p_t$
(the token MLP disabled, as in \eqref{eq:main-anchor-span}) and asks the routed
set states to provide the contextual correction:
\begin{equation}
\label{eq:set-dictionary-span}
f_t=h_t^{(0)}+\operatorname{span}_t,\qquad
\operatorname{span}_t=W_o\!\!\sum_{m\in\mathcal C_t}
\alpha_{t,m}s_m^{(K)}.
\end{equation}
Here $\operatorname{span}_t=W_o r_t$ collects the routed context $r_t$ of
\eqref{eq:main-routed-head} (written here in its single-head form), $s_m^{(K)}$ is
the processed dictionary atom, $W_o$ is the set-output projection
(Section~\ref{sec:model-overview}, identity in the reported runs), and
$\alpha_{t,m}$ is the learned routing weight.  Thus
the contextual part of the token representation must be spanned by causal
dictionary atoms rather than carried by an updated token stream.  This is the
main mathematical distinction from the previous pooling--routing path: pooling
still constructs atoms, and routing still chooses atoms, but the language-model
head no longer receives a full token residual that can silently carry context.

\begin{figure*}[t]
\centering
\includegraphics[width=0.95\textwidth]{set_dictionary_architecture_figure.tikz.pdf}
\caption{Causal dictionary set-attention architecture for the set-dictionary
branch.  The thin vertical rail is the non-contextual anchor
$h_t^{(0)}=e(x_t)+p_t$, carried only to the output residual in
\texttt{anchor\_span} mode.  Fine and coarse set banks are pooled and processed
in parallel, routed back to each token through separate head groups
$\mathcal U_{\mathrm f},\mathcal U_{\mathrm c}$, concatenated, projected to the
span update $\operatorname{span}_t$, and decoded as
$p(x_{t+1}\mid x_{\le t})=\operatorname{softmax}\!\big(\mathbf W_{\mathrm{lm}}(h_t^{(0)}+\operatorname{span}_t)\big)$.
The right panel lists the aligned computation steps used by the
multi-resolution SD-9/SD-9.5 configuration, with fine heads using
$(w_{\mathrm f},s_{\mathrm f})=(2,1)$ and coarse heads using
$(w_{\mathrm c},s_{\mathrm c})=(4,2)$.}
\label{fig:set-dictionary-architecture}
\end{figure*}

The branch also tests two causal candidate fibers.  The endpoint-window fiber
used by the main strict-past experiments is
\begin{equation}
\mathcal C_t^{\mathrm{win}}=\{m:t-w<e_m\le t\},
\end{equation}
where $e_m$ is the endpoint of set $m$.  The maximal causal fiber used in the
\texttt{all\_past} diagnostic is
\begin{equation}
\mathcal C_t^{\mathrm{all}}=\{m:e_m\le t\}.
\end{equation}
Both are causal; \(\mathcal C_t^{\mathrm{all}}\) removes the local
candidate-count ceiling by allowing every sealed past atom to compete for each
token.

Finally, the trained-anchor diagnostic gives the span a predictive target
without using that target at inference.  A shallow causal pre-encoder produces
$h_t^\star$ and has its own next-token CE head.  During training, the objective
is
\begin{equation}
\label{eq:set-dictionary-anchor-loss}
\mathcal L
=\mathcal L_{\mathrm{CE}}
+\lambda_{\mathrm{pre}}\mathcal L_{\mathrm{CE,pre}}
+\lambda_h\,\frac1L\sum_{t=1}^L
\left\|\operatorname{LN}(\operatorname{span}_t)
-\operatorname{sg}\!\left(\operatorname{LN}(h_t^\star)\right)\right\|_2^2,
\end{equation}
where $\operatorname{sg}$ stops gradients through the anchor target.  The
pre-encoder and its head are training-only and are excluded from inference
accounting.  This guard matters: if $h_t^\star$ is not trained, the anchor is a
random target and the reconstruction diagnostic is uninterpretable.

The related experiments are grouped in Table~\ref{tab:sd-capacity-ladder},
Figure~\ref{fig:sd-capacity-ladder}, and
Table~\ref{tab:sd-width-comparison}.  They test three increasingly direct
interventions: a valid trained anchor, the maximal causal fiber, and
over-provisioned dictionary atoms wider than the matched token baseline.

\begin{figure*}[t]
\centering
\includegraphics[width=0.92\textwidth]{../plots/main/fig_sd_capacity_ladder.pdf}
\caption{Set-dictionary capacity ladder at $L=512$ and LR $10^{-4}$.  The
trained-anchor run is valid: the anchor reconstruction error is below the
pre-registered guard threshold and corresponds to a normalized cosine plateau
near $0.30$.  Nevertheless, perplexity is flat.  Widening the candidate fiber
to all sealed past atoms helps the more compressed $(16,8)$ topology but
plateaus far above the dense token baseline; doubling atom width at $(4,2)$
recovers only part of the remaining gap.}
\label{fig:sd-capacity-ladder}
\end{figure*}

\begin{table*}[t]
\centering
\caption{Validated set-dictionary ladder.  All rows use dense exact
set-processing, \texttt{anchor\_span} output, strict-past causality, three
seeds, and WikiText-2 at $L=512$.  The trained-anchor row uses a causal
pre-encoder with its own next-token CE loss; the \texttt{all\_past} rows use
the maximal causal fiber.}
\label{tab:sd-capacity-ladder}
\scriptsize
\resizebox{\textwidth}{!}{%
\begin{tabular}{llccccc}
\toprule
Intervention & Candidate fiber / width & Anchor valid? & Recon. err. & $(16,8)$ PPL $\downarrow$ & $(4,2)$ PPL $\downarrow$ & Span-carried? \\
\midrule
Endpoint CE only & endpoint window, $d_s=d_\phi=384$ & -- & -- & 1510.9 & 1297.9 & yes \\
Trained predictive anchor & endpoint window, $d_s=d_\phi=384$ & yes & 1.18 & 1510.4 & 1289.0 & yes \\
Maximal causal fiber & \texttt{all\_past}, $d_s=d_\phi=384$ & -- & -- & 1363.9 & 1288.6 & yes \\
Over-provisioned atoms & \texttt{all\_past}, $d_s=d_\phi=768$ & -- & -- & -- & 1241.5 & yes \\
\bottomrule
\end{tabular}%
}
\end{table*}

The important pattern is convergence across interventions.  The trained-anchor
run is no longer confounded by an untrained random target: its final
reconstruction error is $1.18$, below the validity guard, and the equivalent
normalized cosine plateau is about $0.30$.  This confirms that a predictive
signal reaches the span, but validation perplexity remains statistically flat
relative to the endpoint CE-only row.  The maximal causal fiber then removes
the local candidate-count ceiling.  It materially improves $(16,8)$ from
$1510.9$ to $1363.9$ PPL, but still remains far from the dense token baseline
at $781.1$ PPL; at $(4,2)$, where the local fiber was already stronger, the
result is essentially unchanged.  Finally, the over-provisioned atom-width
check improves the best set-dictionary $(4,2)$ result to $1241.5$ PPL, but
Table~\ref{tab:sd-width-comparison} shows that this still trails the dense
token baseline by $460.4$ PPL.

\begin{table}[t]
\centering
\caption{Over-provisioned atom-width check at $(w,s)=(4,2)$.  The wide
set-dictionary row uses $d_\phi=d_s=768$, while the matched dense token
baseline has width $D=384$.  VRAM is peak train memory.}
\label{tab:sd-width-comparison}
\scriptsize
\begin{tabular}{lrr}
\toprule
Model & Val. PPL $\downarrow$ & Peak VRAM MiB \\
\midrule
Dense token baseline & 781.1 & 13407.2 \\
Best SD all\_past, $d_s=d_\phi=384$ & 1288.6 & 11913.5 \\
SD all\_past, $d_s=d_\phi=768$ & 1241.5 & 12764.1 \\
\bottomrule
\end{tabular}
\end{table}

These results make the negative result cleaner rather than weaker.  Three
valid capacity probes point to the same ceiling: a trained anchor supplies a
predictive target but the span absorbs only a limited fraction of it; the
maximal causal fiber cannot close the gap; and atoms wider than the token
baseline still remain far behind token attention.  Span-ablation deltas remain
very large throughout, so the anchor is not silently solving the task.  The
most parsimonious explanation is an upstream pooling-stage ceiling: once a bank
atom is formed by compressing token states into a finite set summary, later
routing, anchoring, and atom widening cannot recover all token-resolution
predictive detail.  This is a publishable mechanism result, consistent with
the paper's candidate-count, concentration, and transport framing.

\FloatBarrier

\subsection{Compressive-Memory Evaluation and Compression Baselines}
\label{sec:future-compressive-memory}

The set-dictionary ladder in Section~\ref{sec:future-set-dictionary} suggests
that the next research direction should not be another short-context capacity
sweep.  The better framing is compressive memory: set states should be judged
as learned summaries of past context under explicit memory budgets, not only as
a drop-in replacement for token attention at $L=512$.  A useful follow-up
should therefore pair long-context evaluation with matched compression
baselines.  Examples include sliding-window token attention with the same peak
VRAM, recurrent or chunked memory with the same retained-state budget,
landmark-token baselines with matched memory, and retrieval-style baselines
that expose how much information survives in a compact state.  The key metric
should be a quality--memory frontier over increasing $L$, with per-family
token controls and set-mediated models compared at fixed training budget and
fixed peak memory.

This reframing also clarifies the role of the current negative results.  The
trained-anchor, \texttt{all\_past}, and doubled-width probes show that simply
giving the router more candidates, a predictive target, or wider atoms is not
enough in the short-context token-replacement setting.  But those failures are
precisely the diagnostics one wants before treating the set bank as a
compressed memory.  Future experiments should ask whether the same bank gives a
better long-context tradeoff when the baseline is an explicit compression
system rather than unconstrained dense token attention.  The most informative
protocol is a fixed-VRAM long-context suite: evaluate continuation,
needle-style retrieval, associative recall, and long-document language
modeling while comparing against token-window and compressed-memory baselines
at the same memory budget.

We therefore propose four targeted follow-up experiments after the routing
memory fix.  First, use a compression curriculum: start from nearly
token-equivalent banks such as $(w,s)=(1,1)$ or $(2,1)$ and gradually move to
$(4,2)$ and $(8,4)$, instead of asking the model to learn language modeling,
routing, pooling, and compression simultaneously.  Second, train with a matched
token-attention teacher, adding logit distillation and optionally hidden-state
matching so that set states learn which predictive information must survive
compression.  Third, evaluate under a fixed-VRAM long-context protocol at larger
$L$, comparing dense and landmark token controls against dense and landmark SKA
with five seeds.  Fourth, because Table~\ref{tab:long-context} shows a partial
long-context gap reduction, add an auxiliary compression objective: reconstruct
teacher hidden states, token windows, or next-token distributions from the set
states before or alongside the causal LM loss.

	Two architectural variants are also worth isolating.  The first is activation
	sharing: compute routing queries from a lightweight descriptor rather than
	keeping a full token stream solely for routing, checkpoint or recompute token
	states that are routing-only, or drop token activations after pooling except at
	positions whose candidate fiber is empty.  These variants are more invasive than
	candidate-gather routing and may affect quality, but they address the audit
	observation that the singleton limit carries both token and set states.  The
	second is hybrid attention: use token attention in some layers and set attention
	in others.  This is closer to practical long-context systems, where local/global
	sparse attention, recurrent memory, or retrieval often complements rather than
	fully replaces token attention
	\citep{dai2019transformerxl,beltagy2020longformer,wu2022memorizing}.
	
	\subsection{Hybrid Token--Set Bridge Pilot}
	\label{sec:hybrid-token-set-pilot}
	
	We ran a small hybrid pilot to test whether preserving token-attention layers
	can bridge fine-grained token information through a set-mediated stack.  The
	hybrid model uses one shared token hidden stream $X\in\mathbb R^{B\times L\times
	d}$: token layers and set layers both update this same sequence, rather than
	forming disconnected token and set towers.  A token layer therefore transports
	ordinary token-attention features forward to later layers, while a set layer
	pools the current shared stream into strict-past endpoint sets, processes those
	sets, and routes the resulting context back to token positions.  This shared
	stream is the intended information bridge.  It does not make the architecture
	equivalent to token attention, because positions with nonempty candidate fibers
	still pass through pooling, set attention, routing projections, and the
	\texttt{empty\_only} output policy.
	
	\begin{table*}[t]
	\centering
	\caption{Sparse local-band hybrid pilot at $D=384$, $d_{\mathrm{ff}}=1536$,
	$L=512$, and LR $10^{-4}$.  Hybrid rows use strict-past set layers with
	progressive set topologies $(w,s)=(4,2)$ followed by $(8,4)$.  Sparse token and
	calibrated sparse-set rows are included as operating-point references from the
	validated A7 summary.}
	\label{tab:hybrid-token-set-pilot}
	\scriptsize
	\resizebox{\textwidth}{!}{%
	\begin{tabular}{llrrrrr}
	\toprule
	Model & Pattern / topology & Seeds & Val. PPL $\downarrow$ & Train PPL & Peak VRAM MiB & Comment \\
	\midrule
	Token sparse & -- & 5 & 742.4 & 228.7 & 13408.7 & Matched sparse local-band baseline. \\
	Set sparse \texttt{empty\_only} & $(1,1)$ & 5 & 747.9 & 228.8 & 13936.6 & Near-token calibration point. \\
	Set sparse \texttt{empty\_only} & $(4,2)$ & 3 & 1438.5 & 481.3 & 11807.3 & First near-2 compressed set point. \\
	\midrule
	Hybrid sparse & TTSSSS, $4{:}2;4{:}2;8{:}4;8{:}4$ & 3 & 2623.9 & 946.3 & 13407.1 & Token bridge only in first two layers. \\
	Hybrid sparse & TSTSTS, $4{:}2;4{:}2;8{:}4$ & 3 & 2576.9 & 911.8 & 13517.2 & Alternating token/set bridge. \\
	Hybrid sparse & TTTTSS, $4{:}2;8{:}4$ & 3 & 2393.4 & 881.5 & 13401.1 & Best pilot; still far from controls. \\
	\bottomrule
	\end{tabular}%
	}
	\end{table*}
	
	Table~\ref{tab:hybrid-token-set-pilot} shows a clear negative result, but the
	diagnostics make it informative.  Keeping more token layers helps monotonically
	among the tested patterns: TTTTSS is better than TSTSTS and TTSSSS, which
	supports the bridge intuition.  However, even the best hybrid row remains much
	worse than both the sparse token baseline and the single-stack sparse set model
	at comparable compression.  This is not simply a candidate-count failure.  The
	single-stack sparse set row at $(w,s)=(4,2)$ has mean candidate count about
	$1.98$ and reaches PPL $1438.5$, whereas the hybrid TTTTSS uses the same
	near-2 set topology in one set layer plus an $(8,4)$ layer with mean candidate
	count about $1.96$, yet degrades to PPL $2393.4$.  Repeated compression through
	multiple set layers appears more damaging than one compressed set stack, even
	when token layers provide a shared-stream bridge.
	
	The broader architecture evidence points in the same direction.  Capacity
	ablations in Table~\ref{tab:capacity-diagnostics} show that widening $d_\phi$,
	widening set states, or increasing set-stack depth does not reliably close the
	gap; in fact, depth 8 and 10 worsen validation perplexity.  The A7
	quality--efficiency frontier shows that compressed set rows can reduce VRAM but
	pay large PPL gaps, while the A8 $L=8192$ linear follow-up shows a strong memory
	benefit without a matching quality recovery.  The hybrid pilot adds one more
	constraint: token/set layer mixing alone does not solve either side of the
	tradeoff.  It also does not create a memory win in this implementation: peak
	VRAM is essentially the same as sparse token attention because the model still
	carries the token stream, set states, and dense token-to-set routing tensors.
	
	The minimal promising path is therefore not ``add set layers to a token model''
	by itself.  It should combine hybrid placement with changes that target the
	measured failure modes.  First, candidate-gather routing should remove dense
	$[B,H,L,M]$ routing tensors; this is mainly a memory fix, but it restores the
	purpose of compression.  Second, the set layers need a training signal that
	tells them what token-level information must survive compression, such as a
	compression curriculum, logit distillation, or hidden-state reconstruction from
	a matched token teacher.  Third, if hybrid layers remain useful, they should
	include explicit residual or cross-layer bridges from early token states into
	later set layers, so compression is not forced to preserve all predictive detail
	through a single routed context.  The observed ordering
	TTTTSS $<$ TSTSTS $<$ TTSSSS in perplexity suggests that a still milder
	pattern, such as a late single set layer or a gated set residual added on top of
	token attention, is a more plausible next test than deeper set-heavy stacks.
	
\FloatBarrier

\subsection{Long-Context Data and Memory Tasks}
\label{sec:future-long-context-memory-tasks}

Finally, the larger-data question should be tested, but not treated as the first
fix.  The 30-epoch convergence panel argues against ``just train longer'' as the
whole explanation, yet the set path has additional routing and compression
machinery that may be undertrained on WikiText-2 scale.  A serious follow-up is
a fixed-token-budget run on WikiText-103 or a fixed OpenWebText subset, using a
compact grid of dense, sparse, and landmark token baselines against matching SKA
families at $(w,s)\in\{(1,1),(2,1),(4,2),(6,2),(16,8)\}$.  This would test whether the
observed compression loss is fundamental under the current architecture or a
small-data/training-regime artifact.

Set states should also be evaluated directly as compressed memories, not
only through next-token perplexity.  Memory and retrieval continuation tasks are
a natural fit: books, code, math, long-document continuation, associative
recall, and needle-style retrieval can test whether a set state preserves the
right information under compression.  This follows the motivation behind
compressive and memorizing sequence models, where useful long-context behavior
comes from storing or retrieving compact representations of past inputs rather
than from exact short-context token replacement
\citep{rae2019compressive,wu2022memorizing}.

\FloatBarrier

\subsection{Near-2 Compression Targets}
\label{sec:near2-targets}

The first implementation follow-up should not start with a broad topology
sweep.  The most useful target is the near-2 compression regime, where $L/M$ is
close to 2 and the candidate fiber still has roughly two or more sets per token.
Table~\ref{tab:near2-targets} summarizes the evidence from the causal
empty-only calibration runs.  The important comparison is that $(w,s)=(2,2)$
and $(4,2)$ both have $L/M\simeq 2$, but $(2,2)$ leaves only about one candidate
per token and performs much worse.  Thus compression factor alone is not the
right control variable; near-2 compression should be paired with sufficient
candidate support.

\begin{table*}[t]
\centering
\caption{Near-2 compression evidence from the causal empty-only calibration
summary at $D=384$, $d_{\mathrm{ff}}=1536$, $L=512$, and LR $10^{-4}$.  Values
are mean validation perplexity and peak VRAM over the available seeds from
\texttt{a7\_backend\_family\_empty\_only\_augmented\_summary.tsv}.}
\label{tab:near2-targets}
\scriptsize
\resizebox{\textwidth}{!}{%
\begin{tabular}{llrrrrrl}
\toprule
Family & $(w,s)$ & $M$ & $L/M$ & Mean $|\mathcal C_t|$ & Seeds & Val PPL $\downarrow$ & Interpretation \\
\midrule
Set Dense & $(2,2)$ & 256 & 2.000 & 0.998 & 3 & 1429.0 & Near-2 but too few candidates. \\
Set Dense & $(4,2)$ & 255 & 2.008 & 1.984 & 3 & 1273.6 & Best observed dense near-2 point. \\
Set Dense & $(8,4)$ & 127 & 4.031 & 1.965 & 3 & 1311.7 & More compression, worse PPL. \\
Set Linear & $(2,2)$ & 256 & 2.000 & 0.998 & 3 & 1693.6 & Near-2 but too few candidates. \\
Set Linear & $(4,2)$ & 255 & 2.008 & 1.984 & 3 & 1356.2 & Best observed linear near-2 point. \\
Set Linear & $(8,4)$ & 127 & 4.031 & 1.965 & 3 & 1394.9 & More compression, worse PPL. \\
\bottomrule
\end{tabular}%
}
\end{table*}

This motivates a narrow implementation benchmark after candidate-gather routing:
use $(w,s)=(4,2)$ as the primary quality-preserving compression target and test
$(6,2)$ as the first unobserved variant.  For $L=512$, $(6,2)$ has
$M=\lfloor(512-6)/2\rfloor+1=254$ and $L/M=2.016$, so it preserves the near-2
compression factor while raising the expected interior candidate count from
about 2 to about 3.  The hypothesis is that this additional overlap can reduce
perplexity without leaving the intended compression regime.  The midpoint
$(5,2)$ should be kept as a fallback if $(6,2)$ improves quality but costs too
much memory or over-smooths routing diagnostics.  Larger-$L$ runs should carry
forward whichever of $(4,2)$ and $(6,2)$ wins at $L=512$, starting with
$L=4096$ smoke tests and treating $L=8192$ as a second-stage smoke only after
the matched token and SKA memory paths are validated.

\FloatBarrier
\bibliography{example_paper}
\bibliographystyle{plainnat}

%%%%%%%%%%%%%%%%%%%%%%%%%%%%%%%%%%%%%%%%%%%%%%%%%%%%%%%%%%%%%%%%%%%%%%%%%%%%%%%
%%%%%%%%%%%%%%%%%%%%%%%%%%%%%%%%%%%%%%%%%%%%%%%%%%%%%%%%%%%%%%%%%%%%%%%%%%%%%%%
% APPENDIX
%%%%%%%%%%%%%%%%%%%%%%%%%%%%%%%%%%%%%%%%%%%%%%%%%%%%%%%%%%%%%%%%%%%%%%%%%%%%%%%
%%%%%%%%%%%%%%%%%%%%%%%%%%%%%%%%%%%%%%%%%%%%%%%%%%%%%%%%%%%%%%%%%%%%%%%%%%%%%%%
\newpage
\appendix
\onecolumn
\section{Detailed Model Definition and Structural Objects}
\label{sec:model-objects}

We now formalize the Set Attention language model studied in this work, together with the structural objects used later in the tiered theoretical analysis. The goal of this section is not merely architectural description: it fixes the finite-set, probability, and operator-theoretic objects that appear in Section~\ref{sec:theory}. Every symbol introduced here is used consistently there.

%%%%%%%%%%%%%%%%%%%%%%%%%%%%%%%%%%%%%%%%%%%%%%%%%%%%%%%%%%%%%%%%%%%%%%%%%%%%%%%%
\subsection{Basic Notation}
\label{sec:basic-notation}

\paragraph{Sequence, representation, and capacity parameters.}
Let:
\begin{itemize}
    \item $B \in \mathbb{N}$ denote the batch size,
    \item $L \in \mathbb{N}$ denote the token sequence length,
    \item $L_{\max}\in\mathbb N$ denote the maximum sequence length supported by the positional embedding table,
    \item $D \in \mathbb{N}$ denote the token/state embedding dimension,
    \item $K \in \mathbb N$ denote the number of set-attention blocks,
    \item $H_r \in \mathbb{N}$ denote the number of \emph{routing heads},
    \item $d_r = D/H_r$ denote the dimension per routing head,
    \item $H_a \in \mathbb{N}$ denote the number of \emph{set-attention heads},
    \item $d_a = D/H_a$ denote the dimension per set-attention head,
    \item $d_{\mathrm{ff}} \in \mathbb{N}$ denote the hidden width of the set-level feedforward network,
    \item $V_{\mathrm{vocab}} \in \mathbb{N}$ denote the vocabulary size.
\end{itemize}

\begin{remark}[Router heads vs. set-attention heads]
\label{rem:head-notation}
The routing mechanism and the set-attention stack are distinct architectural modules. Accordingly, we use different symbols for their head counts: $H_r$ for routing heads, and $H_a$ for set-attention heads. This distinction is necessary in Section~\ref{sec:theory}, where routing-head capacity and set-attention stability play different theoretical roles.
\end{remark}

\paragraph{Batch-level vs. single-sample notation.}
Unless stated otherwise, bold uppercase symbols such as $\mathbf{X}^{(0)}$,
$\mathbf{S}^{(0)}$, and $\mathbf{Z}$ denote batched tensors. When the batch
index is suppressed for a fixed sample, we reuse the same base symbol without
boldface, for example $X^{(0)}\in\mathbb R^{L\times D}$ and
$S^{(0)}\in\mathbb R^{M\times D}$. This convention is used throughout
Sections~\ref{sec:model-objects} and~\ref{sec:theory}.

\paragraph{Finite-index and probability notation.}
We write
\[
[L]:=\{1,\dots,L\}
\]
for the token-position index set. For any finite set $A$, let
\[
\Delta(A)
:=
\left\{
p:A\to[0,1]:
\sum_{a\in A} p(a)=1
\right\}
\]
denote the probability simplex on $A$. We use $\delta_a$ for the Dirac mass
concentrated at $a\in A$.
For a linear map $A$, its operator norm is
\[
\|A\|_{\mathrm{op}}:=\sup_{\|v\|_2=1}\|Av\|_2.
\]
We write $\mathbf I_D$ for the $D\times D$ identity matrix.

\paragraph{Bank construction parameters.}
Let:
\begin{itemize}
    \item $w \in \mathbb{N}$ denote the window size (number of tokens per set),
    \item $s \in \mathbb{N}$ denote the stride (spacing between consecutive set starts).
\end{itemize}
We assume throughout that
\begin{equation}
\label{eq:ws-order}
1 \le s \le w \le L.
\end{equation}

\paragraph{Pooling parameters.}
Let:
\begin{itemize}
    \item $\tau_{\mathrm{pool}} > 0$ denote the pooling temperature,
    \item $q \in (0,1)$ denote the trimming quantile,
    \item $\alpha > 0$ denote the trimming sharpness parameter.
\end{itemize}

\paragraph{Routing parameters.}
Let:
\begin{itemize}
    \item $\tau_r > 0$ denote the routing temperature,
    \item $\tau_{\min} > 0$ denote the enforced minimum routing temperature,
    \item $k \in \mathbb{N} \cup \{\infty\}$ denote the optional top-$k$ routing restriction.
\end{itemize}
The implementation clamps the routing temperature below by $\tau_{\min}$; mathematically, this means that all realized routing distributions are evaluated at an effective temperature satisfying
\begin{equation}
\label{eq:tau-clamp}
\tau_r^{\mathrm{eff}} \ge \tau_{\min}.
\end{equation}
In our current experiments, the routing configurations used for the main paper satisfy $H_r>1$, use a fixed finite routing top-$k$ configuration, and pooling is not head-aware. The routing top-$k$ value was not swept as an experimental control.

\paragraph{Embedding matrices.}
The token embedding matrix is
\begin{equation}
\label{eq:token-embedding}
\mathbf{E} \in \mathbb{R}^{V_{\mathrm{vocab}} \times D}
\end{equation}
and the positional embedding matrix is
\begin{equation}
\label{eq:pos-embedding}
\mathbf{P} \in \mathbb{R}^{L_{\max} \times D},
\end{equation}
where positional embeddings are learned absolute embeddings.

\paragraph{Input token sequence.}
An input token sequence is
\begin{equation}
\label{eq:input-sequence}
\mathbf{x}=(x_1,\dots,x_L)\in\{1,\dots,V_{\mathrm{vocab}}\}^L.
\end{equation}

\paragraph{Initial token states.}
For a fixed sample, the embedded token representation is
\begin{equation}
\label{eq:sample-token-state}
h_t^{(0)}=e(x_t)+p_t,
\qquad t\in[L],
\end{equation}
where $e(\cdot)$ denotes the learned token embedding and $p_t$ the learned positional embedding. Batched token representations collect these states as
\begin{equation}
\label{eq:token-states}
\mathbf{X}^{(0)}\in\mathbb{R}^{B\times L\times D},
\end{equation}
with $\mathbf X^{(0)}_{b,t,:}=h_{b,t}^{(0)}$.

\paragraph{Loss and perplexity.}
The training objective is the standard autoregressive cross-entropy loss. For a batch, let
\begin{equation}
\label{eq:loss}
\mathcal{L}\in\mathbb{R}
\end{equation}
denote the batch-averaged loss, and define validation perplexity by
\begin{equation}
\label{eq:perplexity}
\operatorname{PPL}=\exp(\mathcal{L}_{\mathrm{val}}).
\end{equation}

%%%%%%%%%%%%%%%%%%%%%%%%%%%%%%%%%%%%%%%%%%%%%%%%%%%%%%%%%%%%%%%%%%%%%%%%%%%%%%%%
\subsection{Bank Topology as a Hypergraph Incidence Structure}
\label{sec:bank-topology}

The set bank is a deterministic family of overlapping local subsets of token positions. This family is the primary finite-set object underlying the model. For the autoregressive paper mode, we use the strict-past construction: partial trailing windows are dropped, and token-to-set candidates are determined by set endpoints rather than by membership.

\paragraph{Number of sets.}
Given $L$, $w$, and $s$, the number of sets is
\begin{equation}
\label{eq:num-sets}
M=\left\lfloor\frac{L-w}{s}\right\rfloor+1.
\end{equation}

\paragraph{Set family.}
For each set index $m\in\{1,\dots,M\}$, define its start position
\begin{equation}
\label{eq:set-start}
a_m=1+(m-1)s
\end{equation}
and endpoint
\begin{equation}
\label{eq:set-endpoint}
e_m=a_m+w-1.
\end{equation}
The corresponding active token subset is
\begin{equation}
\label{eq:set-definition}
S_m=\{\,t\in[L]: a_m\le t\le e_m\,\}\subseteq[L].
\end{equation}
Thus the bank is the finite set family
\begin{equation}
\label{eq:set-family}
\mathcal{S}=\{S_1,\dots,S_M\}\subseteq 2^{[L]}.
\end{equation}

\paragraph{Trailing tokens and early positions.}
Dropping partial trailing windows means that every active set has cardinality $w$. For arbitrary $(L,w,s)$, membership in the active set family need not cover a final suffix if the last endpoint is less than $L$. Conversely, under strict-past endpoint routing, early token positions can have no available set endpoint. In that case the routed context is zero and the direct token residual carries the position.

\begin{definition}[Bank topology]
\label{def:bank-topology}
The \emph{bank topology} induced by $(L,w,s)$ is the finite hypergraph
\begin{equation}
\label{eq:bank-topology}
\mathcal{H}_{L,w,s}=([L],\mathcal S),
\qquad
[L]=\{1,\dots,L\},
\end{equation}
whose vertex set is the token-position index set and whose hyperedges are the sets in $\mathcal S$.

Equivalently, this hypergraph is represented by its incidence matrix
\begin{equation}
\label{eq:incidence-matrix}
\mathbf A\in\{0,1\}^{L\times M},
\end{equation}
with entries
\begin{equation}
\label{eq:incidence-entry}
A_{t,m}=\mathbbm{1}_{S_m}(t)=
\begin{cases}
1,& t\in S_m,\\
0,& t\notin S_m.
\end{cases}
\end{equation}
\end{definition}

\begin{remark}[Meaning of ``topology'']
\label{rem:topology-meaning}
In this paper, ``topology'' does not refer to a topological space in the sense of open sets. It refers instead to the finite bank structure $\mathcal H_{L,w,s}$ together with the active candidate rule. In strict-past mode, routing support is determined by endpoints; in the named noncausal mode, it is determined by membership. This structure determines routing support, candidate multiplicity, and the topology-induced feasible region analyzed in Section~\ref{sec:theory}.
\end{remark}

\begin{definition}[Candidate family and candidate count]
\label{def:candidate-family}
For each token position $t\in [L]$, define its strict-past \emph{candidate family}
\begin{equation}
\label{eq:candidate-family}
\mathcal C_t=\{m\in\{1,\dots,M\}: t-w<e_m\le t\},
\end{equation}
and its \emph{candidate count}
\begin{equation}
\label{eq:candidate-count}
C_t=|\mathcal C_t|.
\end{equation}
The mean candidate count is
\begin{equation}
\label{eq:mean-candidate-count}
\bar C=\frac{1}{L}\sum_{t=1}^L C_t.
\end{equation}
\end{definition}

The legacy noncausal mode uses the membership candidate family
\begin{equation}
\label{eq:noncausal-candidate-family}
\mathcal C_t^{\mathrm{nc}}=\{m\in\{1,\dots,M\}:t\in S_m\}=\{m:A_{t,m}=1\}.
\end{equation}
Unless explicitly stated otherwise, the autoregressive experiments and diagnostics in this paper use the strict-past endpoint family \eqref{eq:candidate-family}.

\begin{proposition}[Empty strict-past fibers and singleton token limit]
\label{prop:legacy-empty-strict-past-fibers}
Assume the strict-past endpoint bank above with $L\ge w\ge1$ and endpoints
$e_m=w+(m-1)s$ in 1-indexed notation. Then
\[
C_t=0 \quad\Longleftrightarrow\quad t< w .
\]
Consequently, exactly $w-1$ token positions have an empty strict-past
candidate fiber, and a uniformly sampled token position satisfies
\begin{equation}
\label{eq:empty-fiber-probability}
\mathbb P(C_t=0)=\frac{w-1}{L}.
\end{equation}
In the singleton calibration limit $w=1,s=1$, we have $M=L$, $C_t=1$ for
every $t\in[L]$, and therefore $\mathbb P(C_t=0)=0$.
\end{proposition}

\begin{proof}
The earliest sealed endpoint is $e_1=w$. If $t<w$, no endpoint can satisfy
$e_m\le t$, hence $\mathcal C_t=\varnothing$. If $t\ge w$, endpoint $e_1=w$
satisfies $t-w<w\le t$, so $1\in\mathcal C_t$ and $C_t>0$. This proves the
equivalence and the count of empty positions. Dividing by $L$ gives
\eqref{eq:empty-fiber-probability}. When $w=1,s=1$, the starts are
$a_m=m$, the endpoints are $e_m=m$, and $M=L$. For each token $t$, the
candidate condition becomes $t-1<e_m\le t$, which selects exactly
$e_m=t$.
\end{proof}

\paragraph{Residual policy for calibration.}
Proposition~\ref{prop:legacy-empty-strict-past-fibers} motivates separating the
ordinary residual architecture from calibration experiments. The default
\texttt{direct} policy uses
\[
\mathbf z_t=\mathbf X^{(0)}_{t,:}+\mathbf r_t,
\]
which is the architecture used by the main causal runs. The calibrated
\texttt{empty\_only} policy uses
\[
\mathbf z_t=
\begin{cases}
\mathbf X^{(0)}_{t,:},& C_t=0,\\
\mathbf r_t,& C_t>0.
\end{cases}
\]
Thus early strict-past tokens with no sealed endpoint still have a valid
language-model state, while tokens with nonempty set support are evaluated
through the set-routed path without an extra direct token skip. In the
$w=1,s=1$ singleton limit, the empty case never occurs, so
\texttt{empty\_only} coincides with the no-final-residual routed policy.
This makes it the appropriate named mode for empirical convergence tests as
$M/L\to1$ without introducing a residual advantage that is absent from the
matched token-attention baseline.

\begin{proposition}[Interior candidate count]
\label{prop:legacy-interior-candidate-count}
For strict-past interior token positions $t$ sufficiently far from sequence boundaries and after warmup, $C_t$ is the number of stride-spaced endpoints in the interval $(t-w,t]$.
Consequently, for sequences with $L\gg w$,
\begin{equation}
\label{eq:mean-candidate-approx}
\bar C\approx\frac{w}{s}.
\end{equation}
\end{proposition}

\begin{proof}
A set endpoint contributes to the strict-past candidate family if and only if
\[
t-w<e_m\le t.
\]
Substituting $e_m=w+(m-1)s$ gives a length-$w$ interval of admissible set endpoints on a stride-$s$ grid:
\[
\frac{t-w-w}{s}< m-1\le \frac{t-w}{s}.
\]
Thus the exact count can differ by one depending on integer alignment, but is on the order of $w/s$. Averaging over tokens and noting that warmup and boundary effects are negligible when $L\gg w$ gives $\bar C\approx w/s$.
\end{proof}

%%%%%%%%%%%%%%%%%%%%%%%%%%%%%%%%%%%%%%%%%%%%%%%%%%%%%%%%%%%%%%%%%%%%%%%%%%%%%%%%
\subsection{Pooling Operator as Probability Aggregation on Finite Sets}
\label{sec:pooling-operator}

Each set $S_m$ is aggregated into a single embedding by constructing a probability measure on the finite set $S_m$ and taking the expectation of the token-state function under that measure.

\paragraph{Token-state function on a set.}
For each batch element (suppressed when unambiguous) and each set $S_m$, define the token-state function
\begin{equation}
\label{eq:set-token-function}
\mathbf X^{(m)}:S_m\to\mathbb R^D,
\qquad
t\mapsto \mathbf X_{t,:}.
\end{equation}

\paragraph{Uniform set mean.}
Define the set mean embedding by
\begin{equation}
\label{eq:set-mean}
\boldsymbol\mu_m=\frac{1}{|S_m|}\sum_{t\in S_m}\mathbf X_{t,:}
\in\mathbb R^D.
\end{equation}

\paragraph{Normalized squared deviations.}
For each $t\in S_m$, define
\begin{equation}
\label{eq:normalized-deviation}
\delta_{t,m}=\frac{\|\mathbf X_{t,:}-\boldsymbol\mu_m\|_2^2}{D}\in\mathbb R_{\ge 0}.
\end{equation}
These deviations define the energy-like quantities used in the Boltzmann pooling rule.

\paragraph{Quantile threshold and soft trimming gate.}
Let
\begin{equation}
\label{eq:trimming-threshold}
\theta_m=\operatorname{Quantile}_q\bigl(\{\delta_{t,m}:t\in S_m\}\bigr).
\end{equation}
During backpropagation, $\theta_m$ is treated as gradient-detached. Define the sigmoid gate
\begin{equation}
\label{eq:sigmoid}
\sigma(u)=\frac{1}{1+e^{-u}},
\end{equation}
and the soft trimming factor
\begin{equation}
\label{eq:trimming-gate}
r^{\mathrm{trim}}_{t,m}=\sigma\bigl(\alpha(\theta_m-\delta_{t,m})\bigr)\in(0,1).
\end{equation}

\paragraph{Pooling logits and probability mass function.}
Define the pooling logits by
\begin{equation}
\label{eq:pooling-logits}
\ell^{\mathrm{pool}}_{t,m}=-\frac{\delta_{t,m}}{\tau_{\mathrm{pool}}}+\log(r^{\mathrm{trim}}_{t,m}+\varepsilon_0),
\qquad
\varepsilon_0=10^{-8}.
\end{equation}
Then define the pooling weights
\begin{equation}
\label{eq:pooling-weights}
\omega_{t,m}=\frac{\exp(\ell^{\mathrm{pool}}_{t,m})}{\sum_{u\in S_m}\exp(\ell^{\mathrm{pool}}_{u,m})},
\qquad t\in S_m.
\end{equation}
These weights define a probability measure
\begin{equation}
\label{eq:pooling-pmf}
\mathbb P_m\in\Delta(S_m),
\qquad
\mathbb P_m(\{t\})=\omega_{t,m}.
\end{equation}

\begin{remark}[Strict support vs. effective support]
\label{rem:support-vs-effective-support}
Because the implemented pooling rule uses a smooth softmax and a smooth sigmoid gate, the strict support of $\mathbb P_m$ is typically all of $S_m$. For this reason, the more informative notion of concentration is not strict support but \emph{effective support}, introduced later in Section~\ref{sec:theory} through the inverse participation ratio.
\end{remark}

\paragraph{Pooled set embedding as expectation.}
The pooled embedding for set $S_m$ is the expectation of the token-state function under $\mathbb P_m$:
\begin{equation}
\label{eq:pooled-embedding}
s_m^{(0)}
=
\mathbb E_{\mathbb P_m}[\mathbf X^{(m)}]
=
\sum_{t\in S_m}\omega_{t,m}\mathbf X_{t,:}
\in\mathbb R^D.
\end{equation}

\paragraph{Set-state tensor.}
Collecting all pooled set embeddings yields the set-state tensor
\begin{equation}
\label{eq:set-states}
\mathbf S^{(0)}\in\mathbb R^{B\times M\times D},
\qquad
\mathbf S^{(0)}_{b,m,:}=s_{b,m}^{(0)}.
\end{equation}
We use $\mathbf S^{(0)}$ as the canonical notation for post-pooling, pre-set-processing states, since this aligns cleanly with the later gradient diagnostics and theoretical objects.

%%%%%%%%%%%%%%%%%%%%%%%%%%%%%%%%%%%%%%%%%%%%%%%%%%%%%%%%%%%%%%%%%%%%%%%%%%%%%%%%
\subsection{Set-Processing Stack}
\label{sec:set-processing}

The pooled set states are processed by a stack of $K$ set-attention blocks.

\paragraph{Set-state recursion.}
Let
\begin{equation}
\label{eq:set-state-levels}
\mathbf S^{(\ell)}\in\mathbb R^{B\times M\times D},
\qquad
\ell\in\{0,1,\dots,K\},
\end{equation}
with $\mathbf S^{(0)}$ from \eqref{eq:set-states}. Each block computes
\begin{align}
\label{eq:block-attention}
\mathbf S^{(\ell+1/2)} &= \mathbf S^{(\ell)} + \operatorname{Dropout}\Bigl(\operatorname{SetAttention}(\operatorname{LayerNorm}(\mathbf S^{(\ell)}))\Bigr),\\
\label{eq:block-ffn}
\mathbf S^{(\ell+1)} &= \mathbf S^{(\ell+1/2)} + \operatorname{Dropout}\Bigl(\operatorname{FFN}(\operatorname{LayerNorm}(\mathbf S^{(\ell+1/2)}))\Bigr).
\end{align}

\paragraph{Causal set mask.}
For autoregressive language modeling, set attention uses a causal mask
\begin{equation}
\label{eq:causal-mask}
\mathbf C\in\{0,1\}^{M\times M},
\qquad
C_{i,j}=
\begin{cases}
1,& i\ge j,\\
0,& i<j.
\end{cases}
\end{equation}
Thus set $i$ may attend only to itself and preceding sets.

\paragraph{Multi-head set attention.}
Within each block, set attention is standard multi-head self-attention with $H_a$ heads. For each attention head $h\in\{1,\dots,H_a\}$,
\begin{align}
\label{eq:qkv-projections}
\mathbf Q^{(h)} &= \mathbf S\mathbf W_Q^{(h)} \in \mathbb R^{M\times d_a},
&\mathbf W_Q^{(h)}\in\mathbb R^{D\times d_a},\\
\mathbf K^{(h)} &= \mathbf S\mathbf W_K^{(h)} \in \mathbb R^{M\times d_a},
&\mathbf W_K^{(h)}\in\mathbb R^{D\times d_a},\\
\mathbf V^{(h)} &= \mathbf S\mathbf W_V^{(h)} \in \mathbb R^{M\times d_a},
&\mathbf W_V^{(h)}\in\mathbb R^{D\times d_a}.
\end{align}
The implemented score matrix is
\begin{equation}
\label{eq:attention-scores}
\widetilde{\mathbf A}^{(h)} \in \mathbb R^{M\times M}
=
\frac{\mathbf Q^{(h)}(\mathbf K^{(h)})^\top}{\sqrt{d_a}}
\;+\;\mathbf G
\;+\;\mathbf B^{(h)}
\;+\;(\mathbf 1_{M\times M}-\mathbf C)(-\infty),
\end{equation}
where $\mathbf G\in\mathbb R^{M\times M}$ denotes the geometry bias, $\mathbf
B^{(h)}\in\mathbb R^{M\times M}$ denotes the per-head content bias, and
$\mathbf 1_{M\times M}$ is the all-ones matrix. The corresponding attention
output is
\begin{equation}
\label{eq:attention-output}
\mathbf O^{(h)}=\operatorname{softmax}(\widetilde{\mathbf A}^{(h)})\mathbf V^{(h)}
\in \mathbb R^{M\times d_a}.
\end{equation}
Concatenation and output projection yield
\begin{equation}
\label{eq:multihead-concat}
\operatorname{SetAttention}(\mathbf S)
=
[\mathbf O^{(1)};\dots;\mathbf O^{(H_a)}]\mathbf W_O,
\qquad
\mathbf W_O\in\mathbb R^{D\times D}.
\end{equation}

\paragraph{Functional set-attention families.}
The implementation supports three functional set-processing families used in the paper bundle. The dense-exact family uses full set-to-set attention,
\begin{equation}
\label{eq:dense-set-attention-family}
\operatorname{Attn}_{\mathrm{dense}}(Q,K,V)=\operatorname{softmax}(A)V,
\end{equation}
with $A$ given by \eqref{eq:attention-scores}. The linear family used in the experiments is the landmark backend: it routes set interactions through a smaller landmark subset $\Lambda\subseteq\{1,\dots,M\}$, giving a compressed surrogate of the form
\begin{equation}
\label{eq:landmark-set-attention-family}
\operatorname{Attn}_{\mathrm{landmark}}(Q,K,V)
=
\operatorname{softmax}(A_{M\Lambda})\,
\operatorname{softmax}(A_{\Lambda M})V.
\end{equation}
The sparse family used in the experiments is local-band attention: each set attends only within a fixed neighborhood in bank-index space,
\begin{equation}
\label{eq:local-band-set-attention-family}
A_{ij}=-\infty\quad\text{whenever}\quad |i-j|>r,
\end{equation}
up to optional global indices and the causal/set masks. We do not include the Nystr\"om backend in the main paper, since it was not part of the functional family emphasized by the reported experiments.

\paragraph{Set-level feedforward network.}
The block feedforward network is
\begin{equation}
\label{eq:ffn}
\operatorname{FFN}(\mathbf s)=\mathbf W_2\,\operatorname{GELU}(\mathbf W_1\mathbf s),
\end{equation}
with
\begin{equation}
\mathbf W_1\in\mathbb R^{d_{\mathrm{ff}}\times D},
\qquad
\mathbf W_2\in\mathbb R^{D\times d_{\mathrm{ff}}}.
\end{equation}

\paragraph{Processed set states and operator notation.}
After $K$ blocks, the processed set states are
\begin{equation}
\label{eq:final-set-states}
\widetilde{\mathbf S}=\mathbf S^{(K)}\in\mathbb R^{B\times M\times D}.
\end{equation}
We write the full set-processing stack as the operator
\begin{equation}
\label{eq:set-stack-operator}
\Phi_{\mathrm{set}}:\mathbb R^{M\times D}\to\mathbb R^{M\times D},
\qquad
\widetilde{\mathbf S}=\Phi_{\mathrm{set}}(\mathbf S^{(0)}).
\end{equation}

%%%%%%%%%%%%%%%%%%%%%%%%%%%%%%%%%%%%%%%%%%%%%%%%%%%%%%%%%%%%%%%%%%%%%%%%%%%%%%%%
\subsection{Multi-Head Routing Operator}
\label{sec:routing-operator}

The routing operator maps processed set states back to token-level representations. This module is multi-head in all main set-only experiments used for the paper.

\paragraph{Head splitting of set states.}
For each processed set state $\widetilde{\mathbf S}_m\in\mathbb R^D$, split its coordinates into $H_r$ routing-head subvectors:
\begin{equation}
\label{eq:head-split}
\widetilde{\mathbf S}_m=\bigl[\widetilde{\mathbf S}_m^{(1)};\dots;\widetilde{\mathbf S}_m^{(H_r)}\bigr],
\qquad
\widetilde{\mathbf S}_m^{(u)}\in\mathbb R^{d_r}.
\end{equation}
%Equivalently,
%\begin{equation}
%\label{eq:head-split-concat}
%\widetilde{\mathbf S}=\operatorname{Concat}(\widetilde{\mathbf S}^{(1)},\dots,\widetilde{\mathbf S}^{(H_r)}),
%\qquad
%\widetilde{\mathbf S}^{(u)}\in\mathbb R^{M\times d_r}.
%\end{equation}

\paragraph{Per-head routing scores.}
For each token position $t$ and routing head $u\in\{1,\dots,H_r\}$, the learned router computes headwise token-query / set-descriptor compatibility scores. Let $d_m$ denote the descriptor used by the router for set $m$. In the hashed-count experimental path used in the paper, the router receives the pre-set-attention pooled state as this descriptor, so $d_m=s_m^{(0)}$; the hashed-count features affect set attention through content bias. The router computes
\begin{equation}
\label{eq:routing-query-key}
q_t^{(u)}=W_{Q,r}^{(u)}h_t^{(0)},
\qquad
k_m^{(u)}=W_{K,r}^{(u)}d_m,
\end{equation}
and raw scores
\begin{equation}
\label{eq:routing-scores}
g^{(u)}_{t,m}
=
\frac{\langle q_t^{(u)},k_m^{(u)}\rangle}{\sqrt{d_\phi}},
\qquad m\in\{1,\dots,M\}.
\end{equation}
Here $d_\phi$ is the router projection dimension; when a single-head router is used, the same expression reduces to the corresponding full-width query--descriptor dot product.

\paragraph{Candidate masking and optional top-$k$ restriction.}
The topology restricts routing to candidate sets in $\mathcal C_t$. The reported experiments keep a fixed finite router top-$k$ setting rather than sweeping it; in the dense-exact operating points emphasized below, this cutoff is effectively inactive whenever $C_t$ is already smaller than the configured value. Define the masked logits
\begin{equation}
\label{eq:masked-logits}
\widetilde g^{(u)}_{t,m}=
\begin{cases}
g^{(u)}_{t,m},& m\in\mathcal C_t \text{ and } m \text{ survives top-}k,\\
-\infty,& \text{otherwise}.
\end{cases}
\end{equation}

\paragraph{Per-head routing probability measures.}
For each token position $t$ and head $u$, define the routing probability measure
\begin{equation}
\label{eq:routing-probability-measure}
\mathbb Q_t^{(u)}\in\Delta(\mathcal C_t)
\end{equation}
when $C_t>0$, with mass function
\begin{equation}
\label{eq:routing-probabilities}
\pi_{t,m}^{(u)}=\mathbb Q_t^{(u)}(\{m\})
=
\frac{\exp(\widetilde g^{(u)}_{t,m}/\tau_r^{\mathrm{eff}})}{\sum_{j=1}^M \exp(\widetilde g^{(u)}_{t,j}/\tau_r^{\mathrm{eff}})}.
\end{equation}
Its support is
\begin{equation}
\label{eq:routing-support}
\operatorname{supp}(\mathbb Q_t^{(u)})
=
\{m\in\mathcal C_t: \pi_{t,m}^{(u)}>0\}
\subseteq \mathcal C_t.
\end{equation}
Each headwise routing distribution satisfies
\begin{equation}
\label{eq:routing-prob-sum}
\sum_{m=1}^M \pi_{t,m}^{(u)}=1.
\end{equation}
When $C_t=0$ in strict-past mode, no softmax is formed for that token and the routed context below is set to zero.

\paragraph{Per-head routed context.}
For each head $u$, define the routed head-subvector as the expectation of the head-split set-state function under $\mathbb Q_t^{(u)}$:
\begin{equation}
\label{eq:per-head-routed}
\mathbf r_t^{(u)}
=
\mathbb E_{\mathbb Q_t^{(u)}}[\widetilde{\mathbf S}^{(u)}]
=
\sum_{m=1}^M \pi_{t,m}^{(u)}\widetilde{\mathbf S}^{(u)}_m
\in\mathbb R^{d_r}.
\end{equation}

\paragraph{Routed context and output residual policy.}
The routed context is the concatenation of headwise routed subvectors:
\begin{equation}
\label{eq:routed-token}
\mathbf r_t
=
\operatorname{Concat}(\mathbf r_t^{(1)},\dots,\mathbf r_t^{(H_r)})
\in\mathbb R^D.
\end{equation}
If $C_t=0$, define $\mathbf r_t=\mathbf 0$. The default final token representation is the direct residual
\begin{equation}
\label{eq:direct-residual-appendix}
\mathbf z_t=\mathbf X^{(0)}_{t,:}+\mathbf r_t.
\end{equation}
For explicitly configured calibration runs, the \texttt{empty\_only} policy
instead uses $\mathbf z_t=\mathbf X^{(0)}_{t,:}$ when $C_t=0$ and
$\mathbf z_t=\mathbf r_t$ when $C_t>0$; the \texttt{none} policy uses
$\mathbf z_t=\mathbf r_t$ throughout strict-past mode.
Collecting all token positions yields the final token-state tensor
\begin{equation}
\label{eq:routed-tokens}
\mathbf Z\in\mathbb R^{B\times L\times D}.
\end{equation}

\paragraph{Head-aggregated routing distribution.}
For analysis and diagnostics, define the head-aggregated routing distribution by
\begin{equation}
\label{eq:aggregated-routing}
\bar\pi_{t,m}=\frac{1}{H_r}\sum_{u=1}^{H_r}\pi_{t,m}^{(u)}.
\end{equation}
This defines a probability measure
\begin{equation}
\label{eq:aggregated-routing-measure}
\bar{\mathbb Q}_t\in\Delta(\mathcal C_t),
\qquad
\bar{\mathbb Q}_t(\{m\})=\bar\pi_{t,m}.
\end{equation}

\paragraph{Operator notation.}
We write the multi-head routing operator as
\begin{equation}
\label{eq:routing-operator}
\mathcal R_{H_r,\tau_r,\mathcal H,k}:\mathbb R^{M\times D}\to\mathbb R^{L\times D},
\qquad
\mathbf R=\mathcal R_{H_r,\tau_r,\mathcal H,k}(\widetilde{\mathbf S}),
\end{equation}
where $\mathcal H=\mathcal H_{L,w,s}$ is the bank topology, and $\mathbf Z$ is formed from $\mathbf X^{(0)}$ and $\mathbf R$ by the configured output residual policy before being passed to the language-model head.

%%%%%%%%%%%%%%%%%%%%%%%%%%%%%%%%%%%%%%%%%%%%%%%%%%%%%%%%%%%%%%%%%%%%%%%%%%%%%%%%
\subsection{Language Model Head and Loss}
\label{sec:lm-head}

The final token representations are mapped to vocabulary logits by a linear readout:
\begin{equation}
\label{eq:output-logits}
\mathbf Y=\mathbf Z\mathbf W_{\mathrm{lm}}^\top+\mathbf b_{\mathrm{lm}}
\in\mathbb R^{B\times L\times V_{\mathrm{vocab}}},
\end{equation}
where
\begin{equation}
\mathbf W_{\mathrm{lm}}\in\mathbb R^{V_{\mathrm{vocab}}\times D},
\qquad
\mathbf b_{\mathrm{lm}}\in\mathbb R^{V_{\mathrm{vocab}}}.
\end{equation}
For a fixed sample, write $y_t\in\mathbb R^{V_{\mathrm{vocab}}}$ for the vocabulary-logit vector at token position $t$. The autoregressive cross-entropy loss functional is
\begin{equation}
\label{eq:cross-entropy}
\mathcal J
=
-\frac{1}{L}\sum_{t=1}^{L}
\log p_{\mathrm{model}}(x^{\mathrm{target}}_t\mid y_t).
\end{equation}
Here the data loader supplies pre-shifted input and target chunks: the model consumes the input chunk and predicts the aligned next-token target chunk at every label position. The empirical batch loss $\mathcal L$ is the average of $\mathcal J$ over all batch and label positions.

%%%%%%%%%%%%%%%%%%%%%%%%%%%%%%%%%%%%%%%%%%%%%%%%%%%%%%%%%%%%%%%%%%%%%%%%%%%%%%%%
\subsection{Complete Architecture Summary}
\label{sec:architecture-summary}

The full Set Attention language model is the composition:
\begin{equation}
\label{eq:full-pipeline}
\mathbf x
\xrightarrow{\mathrm{Embed}}
\mathbf X^{(0)}
\xrightarrow{\mathcal P}
\mathbf S^{(0)}
\xrightarrow{\Phi_{\mathrm{set}}}
\widetilde{\mathbf S}
\xrightarrow{\mathcal R}
\mathbf R
\xrightarrow{+\mathbf X^{(0)}}
\mathbf Z
\xrightarrow{\mathrm{LM\ Head}}
\mathbf Y,
\end{equation}
where:
\begin{itemize}
    \item $\mathcal P=\mathcal P_{w,s,q,\alpha,\tau_{\mathrm{pool}}}$ is the pooling operator on bank sets,
    \item $\Phi_{\mathrm{set}}$ is the set-processing stack,
    \item $\mathcal R=\mathcal R_{H_r,\tau_r,\mathcal H,k}$ is the multi-head routing operator.
\end{itemize}

\paragraph{Architectural parity fact.}
There is no post-routing token-level MLP in the main proposal. The only feedforward network active in the core architecture is the set-level FFN inside each set-attention block. This fact is important in Section~\ref{sec:theory}, where the theoretical analysis focuses on the set-attention and routing mechanism itself rather than on an added post-routing token nonlinearity.

%%%%%%%%%%%%%%%%%%%%%%%%%%%%%%%%%%%%%%%%%%%%%%%%%%%%%%%%%%%%%%%%%%%%%%%%%%%%%%%%
\subsection{Observable Routing Diagnostics}
\label{sec:routing-diagnostics}

This subsection defines the routing diagnostics used later in Section~\ref{sec:theory} and in the experiments.

\paragraph{Per-head routing entropy.}
For each token $t$ and routing head $u$, define the Shannon entropy
\begin{equation}
\label{eq:per-head-entropy-single}
H(\mathbb Q_t^{(u)})
=
-\sum_{m=1}^M \pi_{t,m}^{(u)}\log \pi_{t,m}^{(u)}.
\end{equation}
The mean per-head routing entropy is
\begin{equation}
\label{eq:per-head-entropy-mean}
H_{\mathrm{head}}
=
\frac{1}{LH_r}
\sum_{t=1}^L\sum_{u=1}^{H_r}
H(\mathbb Q_t^{(u)}).
\end{equation}

\paragraph{Aggregated routing entropy.}
For each token position $t$, define the entropy of the aggregated routing distribution
\begin{equation}
\label{eq:aggregated-entropy-single}
H(\bar{\mathbb Q}_t)
=
-\sum_{m=1}^M \bar\pi_{t,m}\log \bar\pi_{t,m}.
\end{equation}
Its mean over tokens is
\begin{equation}
\label{eq:aggregated-entropy}
H_{\mathrm{agg}}
=
\frac{1}{L}\sum_{t=1}^L H(\bar{\mathbb Q}_t).
\end{equation}

\paragraph{Normalized aggregated entropy.}
To align with the sharp structural bounds in Section~\ref{sec:theory}, define the per-token normalized aggregated entropy by
\begin{equation}
\label{eq:normalized-agg-entropy-token}
\overline H_{\mathrm{agg}}(t)
=
\frac{H(\bar{\mathbb Q}_t)}{\log C_t},
\qquad C_t\ge 2,
\end{equation}
and the mean normalized aggregated entropy by
\begin{equation}
\label{eq:normalized-agg-entropy}
\overline H_{\mathrm{agg}}
=
\frac{1}{|\{t:C_t\ge 2\}|}
\sum_{t:C_t\ge 2}
\frac{H(\bar{\mathbb Q}_t)}{\log C_t}.
\end{equation}
This normalization places the quantity in $[0,1]$ and matches the candidate-family-based entropy bounds used later.

\paragraph{Router top-1 concentration.}
The reported router top-1 diagnostic is probability concentration, not modal head agreement. For each token with nonempty candidate set, define
\begin{equation}
\label{eq:router-top1-concentration-token}
\kappa_t=\frac{1}{H_r}\sum_{u=1}^{H_r}\max_{m\in\mathcal C_t}\pi_{t,m}^{(u)}.
\end{equation}
The mean router top-1 concentration is the average of $\kappa_t$ over tokens with $C_t>0$. It lies in $[1/C_t,1]$ per token when all candidates survive masking and equals $1$ for deterministic routing.

\paragraph{Head agreement.}
For each token $t$, define the top-1 agreement score
\begin{equation}
\label{eq:head-agreement-token}
a_t
=
\frac{1}{H_r}
\max_{m\in\mathcal{C}_t}
\sum_{u=1}^{H_r}
\mathbbm{1}_{\left[\arg\max_{j\in\mathcal{C}_t} \pi_{t,j}^{(u)}=m\right]}
\in\left[\frac{1}{H_r},1\right].
\end{equation}
This score is modal top-1 agreement across routing heads; it is not a measure of full distributional agreement.
The mean head agreement is
\begin{equation}
\label{eq:head-agreement}
\mathcal A
=
\frac{1}{L}\sum_{t=1}^L a_t.
\end{equation}

\paragraph{Head diversity via Jensen--Shannon deviation from the mean.}
For probability laws $p,q\in\Delta(A)$ on a common finite support $A$, define
\[
\operatorname{KL}(p\,\|\,q)
:=
\sum_{a\in A} p(a)\log\frac{p(a)}{q(a)}.
\]
For each token $t$ and head $u$, let
\[
m_t^{(u)}=\tfrac12\left(\mathbb Q_t^{(u)}+\bar{\mathbb Q}_t\right).
\]
Define the Jensen--Shannon divergence from the aggregated routing distribution by
\begin{equation}
\label{eq:js-divergence}
\operatorname{JS}(\mathbb Q_t^{(u)}\,\|\,\bar{\mathbb Q}_t)
=
\frac{1}{2}\operatorname{KL}(\mathbb Q_t^{(u)}\,\|\,m_t^{(u)})
+
\frac{1}{2}\operatorname{KL}(\bar{\mathbb Q}_t\,\|\,m_t^{(u)}).
\end{equation}
The mean JS-to-mean diversity is
\begin{equation}
\label{eq:js-diversity}
\mathcal D_{\mathrm{JS}}
=
\frac{1}{LH_r}
\sum_{t=1}^L\sum_{u=1}^{H_r}
\operatorname{JS}(\mathbb Q_t^{(u)}\,\|\,\bar{\mathbb Q}_t).
\end{equation}

%%%%%%%%%%%%%%%%%%%%%%%%%%%%%%%%%%%%%%%%%%%%%%%%%%%%%%%%%%%%%%%%%%%%%%%%%%%%%%%%
\subsection{Gradient Flow Diagnostics}
\label{sec:gradient-flow-diagnostics}

This subsection fixes the gradient-flow observables used later in the mechanistic analysis.

\paragraph{Intermediate states for gradient diagnostics.}
Let:
\begin{itemize}
    \item $\mathbf h\in\mathbb R^{L\times D}$ denote token states immediately before pooling,
    \item $\mathbf z_0\in\mathbb R^{M\times D}$ denote set states immediately after pooling,
    \item $\mathbf z_K\in\mathbb R^{M\times D}$ denote set states immediately after the set-processing stack.
\end{itemize}
These correspond respectively to the tensor stages $\mathbf X^{(0)}$, $\mathbf S^{(0)}$, and $\widetilde{\mathbf S}$ for a fixed batch element, with notation simplified for the gradient analysis.

\paragraph{Pooling gradient-flow ratio.}
Define
\begin{equation}
\label{eq:rho-p}
\rho_p
=
\frac{\|\nabla_{\mathbf z_0}\mathcal L\|_2}{\|\nabla_{\mathbf h}\mathcal L\|_2+\varepsilon_g},
\qquad
\varepsilon_g=10^{-8}.
\end{equation}

\paragraph{Set-processing gradient-flow ratio.}
Define
\begin{equation}
\label{eq:rho-a}
\rho_a
=
\frac{\|\nabla_{\mathbf z_K}\mathcal L\|_2}{\|\nabla_{\mathbf z_0}\mathcal L\|_2+\varepsilon_g}.
\end{equation}

\paragraph{Combined gradient-flow ratio.}
Define
\begin{equation}
\label{eq:rho-pa}
\rho_{pa}
=
\frac{\|\nabla_{\mathbf z_K}\mathcal L\|_2}{\|\nabla_{\mathbf h}\mathcal L\|_2+\varepsilon_g}.
\end{equation}

\paragraph{Exact stabilizer-corrected multiplicative relation.}
These quantities satisfy the exact identity
\begin{equation}
\label{eq:rho-composition-corrected}
\rho_p\rho_a
=
\rho_{pa}
\frac{\|\nabla_{\mathbf z_0}\mathcal L\|_2}{\|\nabla_{\mathbf z_0}\mathcal L\|_2+\varepsilon_g}.
\end{equation}
Hence the idealized multiplicative relation $\rho_{pa}=\rho_p\rho_a$ is recovered in the limit $\varepsilon_g\to 0$.

\paragraph{Interpretation.}
These ratios are used as local diagnostics of gradient transport through the architectural stages:
\begin{itemize}
    \item $\rho_p\ll 1$: pooling-related gradient bottleneck,
    \item $\rho_p\gg 1$: pooling-related gradient amplification,
    \item $\rho_a\ll 1$: set-stack-related gradient decay,
    \item $\rho_a\gg 1$: set-stack-related gradient amplification.
\end{itemize}

%%%%%%%%%%%%%%%%%%%%%%%%%%%%%%%%%%%%%%%%%%%%%%%%%%%%%%%%%%%%%%%%%%%%%%%%%%%%%%%%
% END OF SECTION 4
%%%%%%%%%%%%%%%%%%%%%%%%%%%%%%%%%%%%%%%%%%%%%%%%%%%%%%%%%%%%%%%%%%%%%%%%%%%%%%%%



%%%%%%%%%%%%%%%%%%%%%%%%%%%%%%%%%%%%%%%%%%%%%%%%%%%%%%%%%%%%%%%%%%%%%%%%%%%%%%%%
% SECTION 5: TIERED THEORETICAL ANALYSIS
%%%%%%%%%%%%%%%%%%%%%%%%%%%%%%%%%%%%%%%%%%%%%%%%%%%%%%%%%%%%%%%%%%%%%%%%%%%%%%%%

%%%%%%%%%%%%%%%%%%%%%%%%%%%%%%%%%%%%%%%%%%%%%%%%%%%%%%%%%%%%%%%%%%%%%%%%%%%%%%%%
% SECTION 5: TIERED THEORETICAL ANALYSIS
%%%%%%%%%%%%%%%%%%%%%%%%%%%%%%%%%%%%%%%%%%%%%%%%%%%%%%%%%%%%%%%%%%%%%%%%%%%%%%%%



\section{Current Multiresolution Theory and Proofs}
\label{sec:theory}

This appendix is the canonical theory for the current set-dictionary branch.  It supersedes the older single-resolution \texttt{direct}/\texttt{empty\_only} proof sketch for the claims cited from Section~\ref{sec:main-theory}.  The executed branch uses exact dense set attention, strict-past banks, endpoint-window candidate gather, fine groups $(w,s,H)=(2,1,6)$, coarse groups $(4,2,2)$, $D=384$, $d_{\rm ff}=1536$, $K=6$, $H=8$, token MLP disabled, trained anchor disabled, hashed-count features, and \texttt{anchor\_span} readout.

\subsection{Notation and Dimensions}
\label{sec:mrp6-notation}

Let $x_0,\ldots,x_{L-1}$ be a token sequence.  A resolution group $g\in\mathcal G$ has window $w_g$, stride $s_g$, routing-head set $\mathcal U_g$, head count $H_g=|\mathcal U_g|$, and stream width $D_g=D_{\rm set}H_g/H$, with $\sum_gH_g=H$.  The implementation enforces divisibility so each head slice has width $D_g/H_g$.  In the active rows $D_{\rm set}=D=384$, so concatenated routed streams already have model width; $W_o$ is kept explicit for code paths with $D_{\rm set}\ne D$.

The thin anchor is
\begin{equation}
h_t^0=e(x_t)+p_t.
\end{equation}
For group $g$, define
\begin{equation}
M_g(L)=\begin{cases}
0,&L<w_g,\\
\lfloor(L-w_g)/s_g\rfloor+1,&L\ge w_g,
\end{cases}
\qquad
a_m^g=ms_g,
\qquad
e_m^g=a_m^g+w_g-1,
\end{equation}
for $m=0,\ldots,M_g-1$.  The group set is $S_m^g=\{a_m^g,\ldots,e_m^g\}$ and the endpoint-window fiber is
\begin{equation}
\mathcal C_t^g=\{m:t-w_g<e_m^g\le t\},\qquad c_t^g=|\mathcal C_t^g|.
\end{equation}

\subsection{Bank Construction}
\label{sec:mrp6-bank}

\begin{lemma}[Group-local empty and interior fibers]
\label{prop:empty-strict-past-fibers}
For group $g$, $\mathcal C_t^g$ is empty exactly when there is no integer $m$ such that
\begin{equation}
0\le m\le M_g-1,
\qquad
\left\lceil\frac{t-2w_g+2}{s_g}\right\rceil\le m\le \left\lfloor\frac{t-w_g+1}{s_g}\right\rfloor.
\end{equation}
Equivalently, with
\begin{equation}
\ell_t^g=\max\!\left(0,\left\lceil\frac{t-2w_g+2}{s_g}\right\rceil\right),\qquad
\nu_t^g=\min\!\left(M_g-1,\left\lfloor\frac{t-w_g+1}{s_g}\right\rfloor\right),
\end{equation}
we have $c_t^g=\max(0,\nu_t^g-\ell_t^g+1)$.  If $s_g\le w_g$ and boundary truncations are inactive, $c_t^g$ is the number of stride-grid endpoints in $(t-w_g,t]$, hence is either $\lfloor w_g/s_g\rfloor$ or $\lceil w_g/s_g\rceil$.  For the active fine and coarse groups the interior count is $2$.
\end{lemma}

\begin{proof}
Substituting $e_m^g=ms_g+w_g-1$ into $t-w_g<e_m^g\le t$ gives
$t-2w_g+1<ms_g\le t-w_g+1$, which is the displayed integer interval after applying ceiling and floor and intersecting with valid bank indices.  The count of integers in a nonempty closed integer interval is upper endpoint minus lower endpoint plus one.  The endpoint-local statement follows because every candidate endpoint lies in the integer interval $(t-w_g,t]$, and a length-$w_g$ interval intersects an $s_g$-spaced grid in either floor or ceiling many points.  For $(2,1)$ and $(4,2)$ this value is exactly $2$ away from boundaries.  This proof uses the moving recent endpoint grid and does not keep the first endpoint in all later fibers.
\end{proof}

\begin{proposition}[Interior candidate count]
\label{prop:interior-candidate-count}
Under the assumptions of Lemma~\ref{prop:empty-strict-past-fibers}, interior token positions have group-local candidate count
\begin{equation}
c_t^g=\left\lfloor\frac{t-w_g+1}{s_g}\right\rfloor-
\left\lceil\frac{t-2w_g+2}{s_g}\right\rceil+1.
\end{equation}
The count is a structural property of the group bank and does not by itself imply a quality ordering.
\end{proposition}

\begin{proof}
This is Lemma~\ref{prop:empty-strict-past-fibers} with the boundary clamps removed.  The final sentence is a scope statement: quality depends on learned atoms, routing, optimization, and data, none of which appears in the counting formula.
\end{proof}

\subsection{Group-Local Pooling, Features, Set Stacks, and Routing}
\label{sec:mrp6-group-local}

Each group pools the thin anchors inside its own bank:
\begin{equation}
z_m^{0,g}=\mathcal P_g(\{h_j^0:j\in S_m^g\}).
\end{equation}
The active pooling rule is soft-trimmed Boltzmann pooling.  Its pure Boltzmann factor has weights proportional to $\exp(-\delta_{j,m}/\tau_{\rm pool})$; the implementation multiplies that factor by a differentiable trimming gate and renormalizes.  Hashed-count and geometry descriptors are group-local deterministic functions of tokens and positions inside the same bank.  A causal $K$-block set stack $\Phi_g$ produces atoms $z_m^{K,g}$, with the mask $e_j^g\le e_m^g$ for source atom $j$ consumed by atom $m$.

For head $u\in\mathcal U_g$, the router scores only candidates in $\mathcal C_t^g$, applies optional top-$k$ masking, and softmaxes at effective temperature at least $\tau_{\min}$.  If the fiber is empty, the implementation emits the zero route.  Otherwise
\begin{equation}
\pi_{t,m}^{u,g}\ge0,
\qquad
\sum_{m\in\mathcal C_t^g}\pi_{t,m}^{u,g}=1,
\qquad
r_t^{u,g}=\sum_{m\in\mathcal C_t^g}\pi_{t,m}^{u,g}\,\operatorname{slice}_{u,g}(z_m^{K,g}).
\end{equation}
The routed group vector is $r_t^g=\operatorname{Concat}_{u\in\mathcal U_g}r_t^{u,g}$.

\subsection{Concatenated \texttt{anchor\_span} Readout}
\label{sec:mrp6-anchor-span}

The full route and span are
\begin{equation}
r_t=\operatorname{Concat}_{g\in\mathcal G}r_t^g,
\qquad
\operatorname{span}_t=W_or_t.
\end{equation}
The active logits are
\begin{equation}
f_t=h_t^0+\operatorname{span}_t,
\qquad
\ell_t=W_{\rm lm}f_t.
\end{equation}
There is no direct contextual token residual, no token MLP, and no trained anchor in the active branch.  Span ablation sets $r_t:=0$; the remaining readout depends on the current token and position and is not a pure unigram regime.

\subsection{Tier A Exact Results}
\label{sec:mrp6-tier-a}

\begin{theorem}[Causal closure]
\label{thm:mrp6-causal-closure}
For every group $g$, layer $k$, atom $m$, and token $t$, $z_m^{k,g}$ depends only on $x_0,\ldots,x_{e_m^g}$, and $r_t^g$, $r_t$, $\operatorname{span}_t$, and $\ell_t$ depend only on $x_0,\ldots,x_t$.  The model is next-token causal.
\end{theorem}

\begin{proof}
The pooled atom $z_m^{0,g}$ is computed from anchors in $S_m^g$, whose largest index is $e_m^g$; soft-trimmed pooling weights, hashed counts, and geometry descriptors use no later token.  Induct on set-stack depth.  At layer $k+1$, the causal set mask permits atom $m$ to consume only atoms $j$ with $e_j^g\le e_m^g$.  By the induction hypothesis, those source atoms are functions of prefixes contained in $x_{\le e_m^g}$, and layer normalization, score biases, attention, residual updates, and the set FFN are deterministic functions of those inputs.  Thus $z_m^{k+1,g}$ is endpoint-measurable.  At token $t$, routing uses only atoms with $e_m^g\le t$ plus the current anchor query $h_t^0$; empty fibers give a deterministic zero vector.  Concatenation, $W_o$, anchor addition, and $W_{\rm lm}$ preserve dependence on $x_{\le t}$.
\end{proof}

\begin{theorem}[Contextual-path factorization]
\label{thm:contextual-path-factorization}
If two histories $x,x'$ share current token and position at $t$, then
\begin{equation}
\ell_t(x)-\ell_t(x')=W_{\rm lm}W_o\bigl(r_t(x)-r_t(x')\bigr).
\end{equation}
Consequently all historical-context dependence factors through the routed span.  Softmax probabilities cannot distinguish span differences mapped by $W_{\rm lm}W_o$ to an all-classes constant logit shift.
\end{theorem}

\begin{proof}
The readout is $\ell_t(x)=W_{\rm lm}(h_t^0(x)+W_or_t(x))$.  Equal current token and equal position imply $h_t^0(x)=h_t^0(x')$.  Subtracting cancels the anchor and leaves the displayed identity.  The observability qualification is the standard softmax invariance under adding $c\mathbf 1$ to the logits.
\end{proof}

\begin{theorem}[Per-group routing entropy and top-$k$]
\label{thm:per-head-entropy}
For nonempty $\mathcal C_t^g$ and head $u\in\mathcal U_g$,
\begin{equation}
0\le H(\pi_t^{u,g})\le \log |\operatorname{supp}(\pi_t^{u,g})|\le \log c_t^g.
\end{equation}
With top-$k$ masking, $|\operatorname{supp}(\pi_t^{u,g})|\le \min\{c_t^g,k\}$.  If $c_t^g=1$, the group-local routing law is the unique Dirac law; finite logits, learned scores, and temperature cannot change it.  For fixed finite support and fixed masked logits, entropy is nondecreasing in the routing temperature, tends to zero as temperature tends to $0^+$ when the maximum is unique, and tends to the uniform-support entropy as temperature tends to infinity.
\end{theorem}

\begin{proof}
The entropy bound is the finite-support entropy maximum principle, with equality only at the uniform law.  Top-$k$ can leave at most $k$ candidates and cannot add candidates outside $\mathcal C_t^g$.  When $c_t^g=1$, the simplex has one element, so the law is forced.  For the temperature claim, write $p_i(\tau)=\exp(a_i/\tau)/\sum_j\exp(a_j/\tau)$ on a fixed support.  Then $H(p(\tau))=\log Z(\tau)-\tau^{-1}\mathbb E_p[a]$ and differentiation gives $dH/d\tau=\operatorname{Var}_p(a)/\tau^3\ge0$.  The two limits follow from the softmax limits on a finite set.
\end{proof}

\begin{corollary}[Normalized head entropy]
\label{cor:normalized-head-entropy}
If $|\operatorname{supp}(\pi_t^{u,g})|\ge2$, then $H(\pi_t^{u,g})/\log |\operatorname{supp}(\pi_t^{u,g})|\in[0,1]$.
\end{corollary}
\begin{proof}
Divide the preceding entropy bound by the positive denominator.
\end{proof}

\begin{theorem}[Group-local aggregated entropy]
\label{thm:aggregated-vs-head-entropy}
Within a fixed group $g$, the averaged law $\bar\pi_t^g=H_g^{-1}\sum_{u\in\mathcal U_g}\pi_t^{u,g}$ satisfies
\begin{equation}
H(\bar\pi_t^g)\ge H_g^{-1}\sum_{u\in\mathcal U_g}H(\pi_t^{u,g}).
\end{equation}
The statement is not defined across groups unless the atoms are first embedded in the tagged disjoint union.
\end{theorem}
\begin{proof}
Entropy is concave on one finite simplex.  The final sentence records the domain condition: different groups have different atom spaces, so untagged averaging is ill-typed.
\end{proof}

\begin{remark}[Entropy gap interpretation]
\label{rem:entropy-gap-interpretation}
Within a group, the entropy gap measures diversity among heads over the same candidate atom labels.  It is not a cross-resolution diversity measure without explicit tags.
\end{remark}

\begin{proposition}[Support and entropy cap under top-$k$]
\label{prop:topk-support-entropy}
For each group and head, top-$k$ routing gives $|\operatorname{supp}(\pi_t^{u,g})|\le\min\{c_t^g,k\}$ and $H(\pi_t^{u,g})\le\log\min\{c_t^g,k\}$.
\end{proposition}
\begin{proof}
Immediate from the candidate mask, the top-$k$ mask, and Theorem~\ref{thm:per-head-entropy}.
\end{proof}

\begin{theorem}[Group-local degenerate topology]
\label{thm:degenerate-topology}
If $c_t^g=1$ for a group at token $t$, every head in $\mathcal U_g$ routes to the unique candidate, has zero entropy, and is independent of learned scores and temperature.  Empty fibers contribute the zero group route, not an empty-only residual branch.
\end{theorem}
\begin{proof}
The probability simplex over one candidate contains one point.  Empty-fiber behavior is the implementation convention described in Section~\ref{sec:mrp6-group-local}.
\end{proof}

\begin{corollary}[Stride equals window]
\label{cor:stride-equals-window}
If $s_g=w_g$, then every interior token has $c_t^g=1$, so Theorem~\ref{thm:degenerate-topology} applies group-locally.
\end{corollary}
\begin{proof}
Interior endpoint intervals $(t-w_g,t]$ have length $w_g$ and the endpoint grid spacing is exactly $w_g$, so exactly one endpoint lies in the interval.
\end{proof}

\begin{theorem}[Direct-sum routing capacity]
\label{thm:direct-sum-routing-capacity}
For each group, let $\Pi_t^g$ have rows indexed by $u\in\mathcal U_g$ and columns by $\mathcal C_t^g$.  Define $\widehat\Pi_t$ over $\Omega_t=\bigsqcup_g(\{g\}\times\mathcal C_t^g)$ by placing each row only in its group block.  Then
\begin{equation}
\operatorname{rank}(\widehat\Pi_t)=\sum_g\operatorname{rank}(\Pi_t^g)\le\sum_g\min\{H_g,c_t^g\}.
\end{equation}
The admissible assignment space has relative-interior dimension $\sum_g H_g(c_t^g-1)$ over nonempty nonsingleton fibers.
\end{theorem}
\begin{proof}
After permuting rows and columns by group, $\widehat\Pi_t$ is block diagonal with diagonal blocks $\Pi_t^g$.  Rank is additive across block diagonal matrices, and each block rank is bounded by its number of rows and columns.  The dimension formula is the sum of the dimensions of the product of group-head simplices.
\end{proof}

\begin{corollary}[Topology-limited routing capacity]
\label{cor:topology-limited-routing-capacity}
For active interior tokens, the registered ceilings are $2$ for b0, $4$ for b25, $4$ for b50, $4$ for b75, and $2$ for b100.  This distinguishes mixed rows from uniform endpoints but cannot select among b25, b50, and b75.
\end{corollary}
\begin{proof}
Interior fine and coarse fibers both have count $2$.  Substitute the blur head splits into $\sum_g\min\{H_g,2\}$.
\end{proof}

\begin{theorem}[Exact dense score-memory law]
\label{thm:exact-dense-score-memory}
The materialized score tensor count across batch and set-attention layers is
\begin{equation}
A_{\rm score}(L,B)=BK\sum_gH_gM_g(L)^2.
\end{equation}
For fine $(2,1)$ and coarse $(4,2)$, $M_f=L-1$, $M_c=\lfloor(L-4)/2\rfloor+1$, and for even $L$, $M_c=L/2-1$.  With coarse-head fraction $p$,
\begin{equation}
\frac{A_{\rm score}}{BKHL^2}=1-\frac{3p}{4}+O(L^{-1}).
\end{equation}
Replacing one fine head by one coarse head strictly decreases this count for registered $L\ge4$ while preserving $\Theta(L^2)$ exact-dense scaling.
\end{theorem}
\begin{proof}
The exact backend forms a full $B\times H_g\times M_g\times M_g$ score tensor in every group and every layer before applying masks.  Summing these cardinalities over groups and $K$ layers gives the count.  The formulas for $M_f$ and $M_c$ are direct substitutions into the strict-past bank formula.  Since $M_f=L-1$ and $M_c\le L/2-1<M_f$ for $L\ge4$, replacing one fine head changes the count by $BK(M_c^2-M_f^2)<0$.  The asymptotic coefficient follows from $M_g=L/s_g+O(1)$.
\end{proof}

\subsection{Tier B Conditional Results}
\label{sec:mrp6-tier-b}

\begin{theorem}[Routed approximation and span coupling]
\label{thm:mechanistic-decomposition}
For group $g$ and head $u$, suppose implemented candidates are coupled to reference atoms $a_i^{u,g}$ by $\kappa_t^g$, with
\begin{equation}
\|\operatorname{slice}_{u,g}(z_m^{K,g})-a_{\kappa_t^g(m)}^{u,g}\|\le\epsilon_g,
\quad
\|a_i^{u,g}\|\le B_g,
\quad
\operatorname{TV}(\kappa_{t\#}^g\pi_t^{u,g},\rho_t^{u,g})\le\xi_g.
\end{equation}
Then
\begin{equation}
\|r_t^{u,g}-r_{*,t}^{u,g}\|\le\epsilon_g+2B_g\xi_g
\end{equation}
and
\begin{equation}
\|\Delta\operatorname{span}_t\|
\le
\|W_o\|_{\rm op}\left[\sum_gH_g(\epsilon_g+2B_g\xi_g)^2\right]^{1/2}.
\end{equation}
\end{theorem}
\begin{proof}
Let $\mu=\kappa_{t\#}^g\pi_t^{u,g}$.  Insert and subtract the coupled implemented route:
\begin{equation}
r_t^{u,g}-r_{*,t}^{u,g}=\sum_i\mu_i(z_i-a_i)+\sum_i(\mu_i-\rho_i)a_i.
\end{equation}
The first term has norm at most $\epsilon_g$ and the second at most $B_g\|\mu-\rho\|_1=2B_g\operatorname{TV}(\mu,\rho)$.  Concatenating heads gives squared norm bounded by $\sum_gH_g(\epsilon_g+2B_g\xi_g)^2$, and applying $W_o$ gives the span bound.
\end{proof}

\begin{remark}[Mechanistic bound meaning]
\label{rem:mechanistic-bound-meaning}
The bound is an error-propagation result conditional on a chosen reference family, coupling, atom error, transport error, and norm bounds.  It is not a decomposition of perplexity into independent causes.
\end{remark}

\begin{corollary}[Conditional loss stability]
\label{cor:loss-sensitivity-pooling}
If the implemented cross-entropy readout is locally Lipschitz in span with constant $L_{\rm model}$ on a ball containing the implemented and reference spans, then
\begin{equation}
|Q_\theta(h_t^0,\operatorname{span}_t)-Q_\theta(h_t^0,\operatorname{span}_{*,t})|
\le L_{\rm model}\|\Delta\operatorname{span}_t\|.
\end{equation}
Any comparison to a different reference loss additionally requires the reference-family Lipschitz and calibration residual assumptions.
\end{corollary}
\begin{proof}
This is the definition of local Lipschitz continuity applied to the two spans.  A second loss family is controlled only after adding the usual triangle-inequality terms for its own perturbation and the calibration residual at the reference span.
\end{proof}

\begin{theorem}[Pooling collision and recovery boundary]
\label{thm:pooling-collision-recovery}
If two histories with the same current token and position induce identical multigroup processed atoms, descriptors, fibers, and router-query inputs through time $t$, then every deterministic downstream stack/router/readout gives the same span and logits.  Moreover, a continuous pre-stack map from an open subset of $\mathbb R^{LD}$ into $\prod_g\mathbb R^{M_gD_g}$ has no global continuous left inverse on an open set whenever $\sum_gM_gD_g<LD$.
\end{theorem}
\begin{proof}
The first claim is determinism: equal inputs to the downstream maps give equal outputs, and equal current token and position give equal anchors.  For the second claim, a continuous left inverse would make the pre-stack map a continuous injection from an open subset of $\mathbb R^{LD}$ into a lower-dimensional Euclidean space, contradicting invariance of domain and the standard dimension obstruction.
\end{proof}

\begin{theorem}[Discrete interior allocation]
\label{thm:conditional-interior-allocation}
Let $A(n)=(H-n)M_f^2+nM_c^2$ and $E(n)$ be an abstract error for $n$ coarse heads.  If $M_c<M_f$, then $A(n+1)-A(n)<0$.  If $G(n)=E(n-1)-E(n)$ is nonincreasing, $G(1)>0$, and $G(H)<0$, then $E$ has a minimizer in $\{1,\ldots,H-1\}$.
\end{theorem}
\begin{proof}
The memory difference is $M_c^2-M_f^2<0$.  Since $G(1)>0$, $E(1)<E(0)$, so $0$ is not a minimizer.  Since $G(H)<0$, $E(H)>E(H-1)$, so $H$ is not a minimizer.  A finite set has a minimizer, and neither endpoint can be one.
\end{proof}

\subsection{Tier C Empirical Propositions and Conjectures}
\label{sec:mrp6-tier-c}

\begin{proposition}[Effective support identities]
\label{prop:effective-support-bounds}
For any pooling law with weights $\omega_{j,m}^g$ on a group set $S_m^g$,
\begin{equation}
n_{\rm eff}^{m,g}=\frac{1}{\sum_{j\in S_m^g}(\omega_{j,m}^g)^2},
\qquad
1\le n_{\rm eff}^{m,g}\le |S_m^g|.
\end{equation}
\end{proposition}
\begin{proof}
The upper bound follows from Cauchy-Schwarz, $(\sum_j\omega_j)^2\le |S|\sum_j\omega_j^2$, and the lower bound from $\sum_j\omega_j^2\le\sum_j\omega_j=1$.
\end{proof}

\begin{proposition}[Pooling concentration regimes]
\label{prop:pooling-concentration-regimes}
For the pure Boltzmann factor on a fixed finite set, as $\tau_{\rm pool}\to0^+$ the law concentrates on minimum-energy tokens; as $\tau_{\rm pool}\to\infty$ it becomes uniform over finite-energy tokens.  The implemented soft-trimmed rule has the same Boltzmann factor multiplied by a trimming gate, so its high-temperature limit is the normalized trimming gate, not necessarily the full uniform law.
\end{proposition}
\begin{proof}
Both limits are the standard finite softmax limits.  Multiplication by a positive trimming gate changes only the limiting base measure after renormalization.
\end{proof}

\begin{definition}[Current feasible operating region]
\label{def:feasible-operating-region}
A current-model operating region is a fixed control tuple together with group-local nonempty fibers, noncollapsed pooling diagnostics, bounded routed approximation quantities $(\epsilon_g,\xi_g,B_g)$, and local readout stability constants sufficient to make Theorem~\ref{thm:mechanistic-decomposition} and Corollary~\ref{cor:loss-sensitivity-pooling} small.  It is a diagnostic definition, not a universal generalization or perplexity theorem.
\end{definition}

\begin{corollary}[Joint experimental design]
\label{cor:joint-experimental-design}
Empirical specialization must vary or measure bank topology, blur allocation, pooling concentration, routing diagnostics, memory counts, and task loss inside a fixed control tuple.  Rank, entropy, or memory alone cannot select b25 or prove superiority over token attention.
\end{corollary}
\begin{proof}
The exact results control structural support, rank, and memory; the conditional results require approximation, transport, and stability quantities.  None of the structural quantities contains the validation loss function or the data distribution.
\end{proof}

\begin{remark}[What is proven]
\label{rem:what-is-proven}
The appendix proves causal closure, contextual-path factorization, group-local routing entropy, direct-sum routing capacity, exact dense score-memory counts, conditional routed approximation, conditional loss stability, a recovery boundary, and a discrete interior-allocation existence theorem.  It does not prove empirical specialization, perplexity superiority, token-attention equivalence, subquadratic complexity, or an intrinsic maximum feasible context length.
\end{remark}

\subsection{Code and Diagnostic Correspondence}
\label{sec:mrp6-code-correspondence}

\begin{table*}[h]
\centering
\caption{Code-grounded correspondence for the current theory.}
\label{tab:mrp6-code-correspondence}
\small
\begin{tabular}{lll}
\toprule
Theory object & Implementation source & Diagnostic or check \\
\midrule
Strict-past bank, endpoints, fibers & \texttt{src/models/set\_only/banks.py} & candidate-count tests \\
Fine/coarse stream split & \texttt{src/models/set\_only/set\_only\_lm.py} & metadata width and bank tests \\
Causal set stack & \texttt{src/models/set\_only/ska\_block.py}, dense backend mask path & future-token perturbation tests \\
Candidate-gather router & \texttt{src/models/set\_only/router.py} & entropy/top-1 per group \\
\texttt{anchor\_span} readout & \texttt{set\_only\_lm.py} encode path & context-path identity test \\
Dense score tensors & \texttt{src/set\_attention/backends/dense\_exact.py} & score-count formula tests \\
Approximation/allocation arithmetic & theory test files under \texttt{tests/theory/} & seven allocation checks \\
\bottomrule
\end{tabular}
\end{table*}

\subsection{Limitations and Nonclaims}
\label{sec:mrp6-limitations}

The theory uses exact dense score cardinalities, not total VRAM.  Peak VRAM also depends on allocator state, autograd saves, optimizer state, kernel workspaces, and external process occupancy.  OOM observations are fixed hardware/software feasibility outcomes unless separately certified.  The all-fine b0 row is still a set model with $(w,s)=(2,1)$ pooling and routing, not token attention.  The allocation theorem is conditional on marginal error gains and does not identify a blur optimum.  MRP-3-dependent specialization remains conditional until validated mechanism evidence exists.

\subsection{Legacy Label Disposition}
\label{sec:mrp6-legacy-disposition}

\begin{longtable}{p{0.26\textwidth}p{0.14\textwidth}p{0.50\textwidth}}
\caption{Disposition of legacy labels required by MRP-6D.}\label{tab:mrp6-legacy-disposition}\\
\toprule
Legacy label(s) & Status & Current disposition \\
\midrule
\endfirsthead
\toprule
Legacy label(s) & Status & Current disposition \\
\midrule
\endhead
\texttt{prop:empty-strict-past-fibers}, \texttt{prop:interior-candidate-count} & Repaired & Restated per group using the endpoint-grid interval and moving recent endpoints. \\
\texttt{thm:per-head-entropy}, \texttt{cor:normalized-head-entropy} & Retained & Retained per group with finite-temperature and support qualifications. \\
\texttt{thm:aggregated-vs-head-entropy}, \texttt{rem:entropy-gap-interpretation} & Repaired & Retained only inside one group or after explicit tagged-union construction. \\
\texttt{prop:topk-support-entropy} & Repaired & Restated per group with $\min\{c_t^g,k\}$. \\
\texttt{thm:degenerate-topology}, \texttt{cor:stride-equals-window} & Repaired & Group-local collapse; empty fibers are zero routes, not residual policy switches. \\
Common-mixture, strict inclusion, dimension, rank, topology capacity labels & Replaced & Replaced by direct-sum routing capacity and topology-limited capacity over tagged group atoms. \\
\texttt{thm:boltzmann-maxent}, \texttt{rem:soft-trimmed-boltzmann} & Repaired & Pure Gibbs statement requires feasible moments or fixed expected energy; implemented rule is a soft-trimmed reweighting, not the exact max-entropy optimizer. \\
\texttt{thm:gradient-decomposition-pooling}, \texttt{cor:direct-coupling-gradients} & Removed & Not stated for the active full graph because overlapping windows, descriptors, router queries, set-stack paths, and anchor-span readout are active. \\
Raw/logged gradient ratio labels & Partially retained & Algebraic stabilizer identity remains a diagnostic identity; unsupported bottleneck causality is removed. \\
\texttt{thm:necessary-feasible-topology}, \texttt{cor:stride-less-than-window} & Repaired & Structural nondegenerate-routing conditions only, not quality conditions. \\
\texttt{thm:mechanistic-decomposition}, \texttt{rem:mechanistic-bound-meaning} & Repaired & Replaced by multigroup routed approximation and span/loss bounds with explicit reference and Lipschitz assumptions. \\
\texttt{prop:effective-support-bounds}, \texttt{prop:pooling-concentration-regimes} & Retained & Exact probability identities only; direct quality implications removed. \\
\texttt{prop:jacobian-sandwich} & Removed & Invalid for the active graph unless rewritten as branch-isolated VJPs with all active paths named. \\
\texttt{prop:entropy-feasible-routing}, \texttt{prop:capacity-stability} & Repaired & Structural/rank interpretation only; empirical thresholds removed. \\
\texttt{cor:loss-sensitivity-pooling} & Repaired & Replaced by weighted conditional loss-stability statement. \\
\texttt{cor:topology-collapse-boundary} & Repaired & Per-group routing collapse only. \\
\texttt{def:feasible-operating-region}, \texttt{cor:joint-experimental-design}, \texttt{rem:what-is-proven} & Rewritten & Rewritten around exact-dense multiresolution program and explicit nonclaims. \\
Main-text legacy four labels & Replaced & Replaced by contextual-path factorization, direct-sum routing capacity, exact-dense memory coefficient, and conditional interior allocation. \\
\bottomrule
\end{longtable}
\section{Extended Experimental Tables}
\label{app:extended-experimental-tables}

\subsection{Legacy Evidence Exclusion}
\label{app:legacy-evidence-exclusion}

Earlier drafts contained appendix tables and figures generated before the A1
strict-past causal LM correction, or from artifacts whose per-run metadata did
not establish the post-A1 causal LM protocol.  Those displays are excluded from
the active manuscript evidence record.  The underlying files remain preserved
for internal audit only, but reviewer-facing tables, figures, and claims in this
revision use only validated post-A1 causal LM manifests under
\texttt{out/paper\_integrated\_evidence/checks/}.  In particular, legacy
artifact roots such as \texttt{out/paper\_bundle/},
\texttt{out/paper\_complements\_bundle/}, \texttt{out/paper\_complements/},
older \texttt{out/metrics/paper\_action*} outputs, and dense-only mechanism
plots from the pre-A1 period are not mixed with the causal LM results reported
here.

The matched-quality comparison is reported in the main paper as
Table~\ref{tab:matched-baseline-vs-setonly}. Source configurations use the
corresponding implementation names for these paper-level quantities, and the
generated TSV files retain those exact run identifiers.

\iffalse
Legacy pre-A1 / causality-unverified displays intentionally hidden from the
active PDF. They are retained in source only as an audit trail and must not be
used for reviewer-facing causal LM claims unless rebuilt under the post-A1
protocol.

\subsection{Detailed Tier C Tables}
\label{app:tier-c-tables}

The tables in this subsection give rounded numeric values for the Tier~C evidence summarized in the main text. In Table~\ref{tab:anchor-stability}, candidate count is constant at $3.953$ and peak VRAM is constant at $11206.7$ MiB across seeds, so both columns are omitted. Table~\ref{tab:boundary-stride-transition} is retained only as a narrow dense two-point supporting slice; the main paper uses the broader cross-family stride sweep rather than treating this table as standalone evidence for a sharp transition.

\begin{table}[h]
\centering
\caption{Stability at the reference operating point across random seeds.}
\label{tab:anchor-stability}
\scriptsize
\resizebox{\columnwidth}{!}{%
\begin{tabular}{rcccccccc}
\toprule
Seed & Val. PPL $\downarrow$ & Routing entropy & Router top-1 & Pooling $n_{\mathrm{eff}}$ ratio & $\rho_p$ & $\rho_a$ & $\rho_{pa}$ & Time / epoch (s) $\downarrow$ \\
\midrule
0 & 580.6 & 0.264 & 0.395 & 0.411 & 0.829 & 0.128 & 0.106 & 73.4 \\
1 & 653.3 & 0.263 & 0.399 & 0.413 & 0.812 & 0.137 & 0.112 & 99.8 \\
2 & 581.2 & 0.263 & 0.402 & 0.412 & 0.827 & 0.130 & 0.108 & 81.9 \\
Mean $\pm$ std & 605.0 $\pm$ 41.8 & 0.263 $\pm$ 0.001 & 0.399 $\pm$ 0.004 & -- & -- & -- & 0.109 $\pm$ 0.003 & 85.0 $\pm$ 13.5 \\
\bottomrule
\end{tabular}%
}
\end{table}

\begin{table*}[h]
\centering
\caption{Dense topology-temperature slice.}
\label{tab:topology-temperature-feasible-region}
\scriptsize
\resizebox{\textwidth}{!}{%
\begin{tabular}{cccccccccc}
\toprule
$s$ & $w$ & $\tau_r$ & Candidate count & Routing entropy & Router top-1 & $\rho_p$ & $\rho_a$ & $\rho_{pa}$ & Val. PPL $\downarrow$ \\
\midrule
16 & 16 & 0.5 & 1.000 & 0.000 & 1.000 & 1.245 & 0.201 & 0.251 & 1322.1 \\
16 & 16 & 1.0 & 1.000 & 0.000 & 1.000 & 1.224 & 0.201 & 0.247 & 1270.5 \\
16 & 16 & 2.0 & 1.000 & 0.000 & 1.000 & 1.223 & 0.199 & 0.244 & 1362.2 \\
16 & 16 & 4.0 & 1.000 & 0.000 & 1.000 & 1.304 & 0.185 & 0.241 & 1321.3 \\
6 & 16 & 0.5 & 2.641 & 0.195 & 0.537 & 0.944 & 0.146 & 0.138 & 634.6 \\
6 & 16 & 1.0 & 2.641 & 0.199 & 0.518 & 0.934 & 0.151 & 0.141 & 669.7 \\
6 & 16 & 2.0 & 2.641 & 0.201 & 0.507 & 0.921 & 0.152 & 0.140 & 693.7 \\
6 & 16 & 4.0 & 2.641 & 0.203 & 0.496 & 0.914 & 0.151 & 0.138 & 720.3 \\
\bottomrule
\end{tabular}%
}
\end{table*}

\begin{table*}[h]
\centering
\caption{Two-point dense stride/candidate-count check.}
\label{tab:boundary-stride-transition}
\scriptsize
\resizebox{\textwidth}{!}{%
\begin{tabular}{ccccccccccccc}
\toprule
$s$ & $w$ & $\tau_r$ & Candidate count & Routing entropy & Router top-1 & Pooling $n_{\mathrm{eff}}$ ratio & $\rho_p$ & $\rho_a$ & $\rho_{pa}$ & Val. PPL $\downarrow$ & Time / epoch (s) $\downarrow$ & Peak VRAM (MiB) $\downarrow$ \\
\midrule
3 & 16 & 1.0 & 5.266 & 0.300 & 0.332 & 0.413 & 0.786 & 0.128 & 0.101 & 595.4 & 84.5 & 11470.9 \\
5 & 16 & 1.0 & 3.166 & 0.230 & 0.458 & 0.406 & 0.878 & 0.144 & 0.126 & 690.0 & 109.9 & 11064.5 \\
\bottomrule
\end{tabular}%
}
\end{table*}

\begin{table*}[h]
\centering
\caption{Dense pooling-temperature sweep.}
\label{tab:pooling-concentration-transport}
\scriptsize
\resizebox{\textwidth}{!}{%
\begin{tabular}{ccccccccccc}
\toprule
$\tau_{\mathrm{pool}}$ & Pooling $n_{\mathrm{eff}}$ ratio & Pooling $n_{\mathrm{eff}}$ L2 & Pooling effective support & Routing entropy & Router top-1 & $\rho_p$ & $\rho_a$ & $\rho_{pa}$ & Val. PPL $\downarrow$ & Time / epoch (s) $\downarrow$ \\
\midrule
0.05 & 0.209 & 3.277 & 11.156 & 0.262 & 0.399 & 0.432 & 0.232 & 0.100 & 893.4 & 76.3 \\
0.1 & 0.412 & 6.483 & 15.035 & 0.265 & 0.394 & 0.853 & 0.131 & 0.112 & 628.5 & 79.2 \\
0.2 & 0.644 & 10.153 & 15.724 & 0.264 & 0.395 & 1.380 & 0.086 & 0.119 & 389.1 & 94.7 \\
\bottomrule
\end{tabular}%
}
\end{table*}

\clearpage

The matched-quality comparison is reported in the main paper as Table~\ref{tab:matched-baseline-vs-setonly}. Source configurations use the corresponding implementation names for these paper-level quantities, and the generated TSV files retain those exact run identifiers.

\subsection{Cross-Family Complement Sweeps}
\label{app:family-complements}

The following targeted complements use the reference operating recipe with learning rate $10^{-4}$. They expand the two mechanism-critical axes studied in the main paper, stride/candidate-count and pooling temperature, across dense exact, sparse local-band, and linear landmark set-attention families. Rows are deduplicated by the full control key, with audit JSONs retained in the experiment bundle.

\begin{table*}[t]
\centering
\caption{Expanded cross-family boundary-stride complement sweep. Display values are rounded.}
\label{tab:boundary-family-complements}
\small
\begin{tabular}{lccccccccc}
\toprule
Family & $s$ & Cand. & Entropy & Top-1 & $n_{\mathrm{eff}}$ ratio & $\rho_p$ & $\rho_a$ & $\rho_{pa}$ & Val. PPL $\downarrow$ \\
\midrule
Dense & 3 & 5.266 & 0.302 & 0.325 & 0.410 & 0.765 & 0.124 & 0.095 & 570.7 \\
Dense & 4 & 3.953 & 0.265 & 0.393 & 0.410 & 0.874 & 0.134 & 0.117 & 600.5 \\
Dense & 5 & 3.166 & 0.231 & 0.452 & 0.408 & 0.884 & 0.141 & 0.124 & 684.2 \\
Dense & 6 & 2.641 & 0.198 & 0.522 & 0.397 & 0.887 & 0.148 & 0.131 & 733.0 \\
Dense & 8 & 1.984 & 0.152 & 0.625 & 0.403 & 0.975 & 0.165 & 0.161 & 880.1 \\
Linear & 3 & 5.266 & 0.300 & 0.335 & 0.399 & 0.760 & 0.145 & 0.110 & 659.1 \\
Linear & 4 & 3.953 & 0.263 & 0.400 & 0.402 & 0.827 & 0.148 & 0.122 & 706.1 \\
Linear & 5 & 3.166 & 0.231 & 0.455 & 0.402 & 0.888 & 0.151 & 0.134 & 767.1 \\
Linear & 6 & 2.641 & 0.198 & 0.524 & 0.391 & 0.950 & 0.161 & 0.152 & 773.0 \\
Linear & 8 & 1.984 & 0.153 & 0.620 & 0.397 & 1.047 & 0.172 & 0.180 & 1019.1 \\
Sparse & 3 & 5.266 & 0.301 & 0.332 & 0.415 & 0.798 & 0.119 & 0.095 & 549.8 \\
Sparse & 4 & 3.953 & 0.265 & 0.393 & 0.411 & 0.829 & 0.137 & 0.114 & 606.7 \\
Sparse & 5 & 3.166 & 0.231 & 0.453 & 0.411 & 0.864 & 0.144 & 0.125 & 716.4 \\
Sparse & 6 & 2.641 & 0.198 & 0.523 & 0.402 & 0.951 & 0.151 & 0.143 & 664.3 \\
Sparse & 8 & 1.984 & 0.153 & 0.623 & 0.407 & 0.979 & 0.167 & 0.164 & 866.1 \\
\bottomrule
\end{tabular}
\end{table*}

\begin{table*}[t]
\centering
\caption{Expanded cross-family pooling-temperature complement sweep. Display values are rounded.}
\label{tab:pooltau-family-complements}
\small
\begin{tabular}{lccccccccc}
\toprule
Family & $\tau_{\mathrm{pool}}$ & $n_{\mathrm{eff}}$ ratio & $n_{\mathrm{eff}}$ L2 & Eff. support & Entropy & Top-1 & $\rho_{pa}$ & Val. PPL $\downarrow$ & Time / epoch (s) $\downarrow$ \\
\midrule
Dense & 0.050 & 0.204 & 3.200 & 11.063 & 0.260 & 0.406 & 0.098 & 982.5 & 71.3 \\
Dense & 0.075 & 0.313 & 4.922 & 14.060 & 0.264 & 0.395 & 0.103 & 774.1 & 71.3 \\
Dense & 0.100 & 0.411 & 6.466 & 15.034 & 0.263 & 0.401 & 0.110 & 665.0 & 70.6 \\
Dense & 0.150 & 0.557 & 8.774 & 15.591 & 0.265 & 0.393 & 0.120 & 488.4 & 80.9 \\
Dense & 0.200 & 0.642 & 10.114 & 15.714 & 0.264 & 0.394 & 0.125 & 356.9 & 71.3 \\
Linear & 0.050 & 0.207 & 3.233 & 11.035 & 0.263 & 0.398 & 0.098 & 1051.1 & 72.2 \\
Linear & 0.075 & 0.309 & 4.848 & 13.910 & 0.265 & 0.394 & 0.114 & 890.5 & 72.0 \\
Linear & 0.100 & 0.406 & 6.378 & 14.985 & 0.265 & 0.391 & 0.123 & 675.8 & 73.7 \\
Linear & 0.150 & 0.544 & 8.572 & 15.572 & 0.264 & 0.398 & 0.136 & 550.8 & 72.8 \\
Linear & 0.200 & 0.632 & 9.957 & 15.707 & 0.264 & 0.394 & 0.144 & 474.5 & 71.2 \\
Sparse & 0.050 & 0.210 & 3.284 & 11.196 & 0.263 & 0.396 & 0.102 & 872.7 & 71.5 \\
Sparse & 0.075 & 0.314 & 4.917 & 14.020 & 0.264 & 0.397 & 0.110 & 758.5 & 74.9 \\
Sparse & 0.100 & 0.406 & 6.398 & 15.030 & 0.262 & 0.405 & 0.107 & 648.7 & 77.4 \\
Sparse & 0.150 & 0.553 & 8.716 & 15.591 & 0.263 & 0.399 & 0.114 & 462.4 & 72.9 \\
Sparse & 0.200 & 0.641 & 10.104 & 15.725 & 0.263 & 0.398 & 0.119 & 359.6 & 71.9 \\
\bottomrule
\end{tabular}
\end{table*}

\fi

\subsection{LR-Normalized Matched-Grid Provenance}
\label{app:partial-lrnorm-slice}

The completed LR-normalized grids are retained as external provenance artifacts rather than reproduced as an additional appendix table. This avoids mixing older dense-baseline-only slices with the v2.7 matched sparse/linear token controls used for the active claims. Source paths and hashes are recorded in \texttt{audit/A5\_4\_handoff.md} and \texttt{out/paper\_integrated\_evidence/checks/final\_artifact\_index.tsv}; the current paper-level comparison is Table~\ref{tab:matched-baseline-vs-setonly}.

\subsection{Complete Post-A1 Causal LM Experiment Record}
\label{app:complete-experiment-record}

Table~\ref{tab:complete-causal-experiment-record} lists every stabilized
post-A1 causal LM experiment family retained for this revision.  The table is an
index rather than a replacement for the source artifacts: full CSV paths,
manifest paths, row counts, SHA256 hashes, and audit notes are recorded in
\texttt{audit/causal\_experiment\_report.md} and
\texttt{out/paper\_integrated\_evidence/checks/causal\_experiment\_report\_index.tsv}.
This separation keeps the manuscript readable while making the complete
experiment record explicit and recoverable.

\begin{table*}[t]
\centering
\caption{Complete stabilized post-A1 causal LM experiment record used for this
revision.  ``Rows'' is the row count of the corresponding indexed TSV artifact;
all listed manifests passed validation.}
\label{tab:complete-causal-experiment-record}
\scriptsize
\resizebox{\textwidth}{!}{%
\begin{tabular}{llrll}
\toprule
Phase & Experiment family & Rows & Intended paper use & Caveat \\
\midrule
A2 & LR-normalized grid & 153 & Dense baseline and set-family provenance & Matched v2.7 controls supersede backend-family claims \\
A2.4 & Matched token-backend controls & 10 & Sparse/linear token controls at LR-normalized reference & Fixed to $D=384,d_{\mathrm{ff}}=1536,L=512,w=16,s=8$ \\
A3.1 & Fixed-stride window sweep & 18 & Candidate-count and compression mechanism evidence & Set-family mechanism evidence, not headline baseline comparison \\
A3.1-control & Window-sweep token overlays & 12 & Backend attribution control for A3.1 & Use with A3.1 to separate set routing from backend effects \\
A3.2 & Pooling-temperature sweep & 15 & Pooling support and transport diagnostics & High-temperature extension shows backend-specific saturation \\
A3.3 & Stride/candidate-count sweep & 12 & Topology complement for Tier C diagnostics & Demoted complement rather than headline result \\
A4.1 & Long-context smoke & 2 & Feasible long-context batch policy & Not a final comparison table \\
A4.2 & Long-context set-family slice & 12 & Quality/memory comparison at $L=2048$ & Pair with A4.2 controls for backend-matched claims \\
A4.2-control & Long-context token controls & 2 & Matched sparse/linear token controls at $L=2048$ & Dense token baseline remains strongest in perplexity \\
A4.3 & 30-epoch convergence panel & 7 & Convergence check beyond 10 epochs & Token baselines remain favored under tested settings \\
A6.1 & $d_\phi$ interface capacity & 9 & Set-token interface bottleneck diagnostic & Family-specific gains, not a complete fix \\
A6.2 & Set-state width capacity & 12 & Explicit set-state dimensionality test & Wider set states do not reliably close the gap \\
A6.3 & Joint interface bottleneck & 18 & Coupled set-state and $d_\phi$ bottleneck test & Matching $d_\phi$ to set-state width often worsens PPL \\
A6.4 & Set-stack depth sweep & 18 & Depth bottleneck diagnostic & Depth 8/10 worsens validation PPL versus depth 6 \\
A7 & \texttt{empty\_only} calibration & 105 & Empirical token-limit and backend-matched compression paths & Five-seed support at convergence-critical points; convergence as $M/L\to1$, not exact equivalence \\
A8-hybrid & Hybrid token/set sparse pilot & 9 & Token--set bridge diagnostic for mixed layer stacks & Completed cleanly but poor PPL; not worth extending to seeds 3,4 without a new routing/memory fix \\
\bottomrule
\end{tabular}%
}
\end{table*}

\iffalse
Legacy dense-only mechanism figures intentionally hidden from the active PDF.
They were generated before the current post-A1 causal LM evidence boundary and
must not be mixed with the validated causal LM record.

\subsection{Additional Dense-Only Mechanism Views}
\label{app:dense-only-mechanism-figures}

These dense-only figures give an additional focused slice of the same Tier~C mechanisms. The main paper uses the cross-family figures because they test the mechanisms across dense exact, sparse local-band, and linear landmark families. In the topology-temperature views, candidate-count-$1$ settings collapse toward worse validation quality, zero entropy, and unit concentration. In the pooling-temperature views, broader pooling support coincides with lower validation perplexity and stronger end-to-end transport.

\begin{figure*}[t]
\centering
\begin{subfigure}[t]{0.32\textwidth}
\centering
\includegraphics[width=\linewidth]{../plots/main/fig_phase_candidates_vs_ppl.png}
\caption{Validation perplexity versus candidate count.}
\label{fig:phase-candidates-vs-ppl}
\end{subfigure}\hfill
\begin{subfigure}[t]{0.32\textwidth}
\centering
\includegraphics[width=\linewidth]{../plots/main/fig_phase_candidates_vs_entropy.png}
\caption{Routing entropy versus candidate count.}
\label{fig:phase-candidates-vs-entropy}
\end{subfigure}\hfill
\begin{subfigure}[t]{0.32\textwidth}
\centering
\includegraphics[width=\linewidth]{../plots/main/fig_phase_candidates_vs_top1.png}
\caption{Router top-1 versus candidate count.}
\label{fig:phase-candidates-vs-top1}
\end{subfigure}
\caption{Dense-only topology-temperature views.}
\label{fig:dense-only-topology-temperature}
\end{figure*}

\begin{figure*}[t]
\centering
\begin{subfigure}[t]{0.32\textwidth}
\centering
\includegraphics[width=\linewidth]{../plots/main/fig_pooltau_vs_ppl.png}
\caption{Validation perplexity versus pooling temperature.}
\label{fig:pooltau-vs-ppl}
\end{subfigure}\hfill
\begin{subfigure}[t]{0.32\textwidth}
\centering
\includegraphics[width=\linewidth]{../plots/main/fig_pooltau_vs_neff.png}
\caption{Effective-support summaries versus pooling temperature.}
\label{fig:pooltau-vs-neff}
\end{subfigure}\hfill
\begin{subfigure}[t]{0.32\textwidth}
\centering
\includegraphics[width=\linewidth]{../plots/main/fig_pooltau_vs_rho_pa.png}
\caption{End-to-end gradient transport $\rho_{pa}$ versus pooling temperature.}
\label{fig:pooltau-vs-rho-pa}
\end{subfigure}
\caption{Dense-only pooling-temperature views.}
\label{fig:dense-only-pooltau}
\end{figure*}
\fi

\section{Implementation Notes and Supporting Comparisons}
\label{app:implementation-notes}

\subsection{Pre-Pooling Token MLP}
\label{app:token-mlp}

The implementation includes an optional pre-pooling token MLP,
\[
\operatorname{TokenMLP}(h)=W_2\,\operatorname{GELU}(W_1h),
\]
applied after token and positional embeddings and before bank pooling. This module was disabled in the experiments emphasized in the main paper and was not part of the topology, temperature, pooling, dense, sparse, or linear family sweeps. Preliminary checks with the module enabled did not materially change the qualitative behavior of the model, so it is treated as an implementation extension rather than a core architectural component.

\subsection{Dropped Historical Parity Slice}
\label{app:parity-support}

Earlier drafts contained an appendix-only parity slice with baseline values that could not be reconciled with the canonical LR-normalized artifacts. The audit in \texttt{audit/A1\_3\_reconciliation.md} therefore treats that slice as unrecoverable for paper use. The active comparisons in this paper are the matched-control and long-context tables in the main Results section, with provenance recorded in the final reproducibility handoff.
%%%%%%%%%%%%%%%%%%%%%%%%%%%%%%%%%%%%%%%%%%%%%%%%%%%%%%%%%%%%%%%%%%%%%%%%%%%%%%%
%%%%%%%%%%%%%%%%%%%%%%%%%%%%%%%%%%%%%%%%%%%%%%%%%%%%%%%%%%%%%%%%%%%%%%%%%%%%%%%

\input{neurips_checklist}

\end{document}
